%% Empirical Software Engineering (Springer) -- SVJour3.
%%
%% GENERATED by make_emse.py from manuscript3_revised.tex. Edit the source and
%% regenerate; edits made here are lost on the next build.
%%
%% Build with pdfLaTeX. svjour3.cls and svglov3.clo are in this folder.
%% The bibliography is an inline thebibliography, so BibTeX never runs and no
%% .bst is required; spbasic.bst is included only in case you later move to a
%% .bib file.
\RequirePackage{fix-cm}
\documentclass[smallextended,natbib]{svjour3}

\smartqed

\usepackage{amsmath}
\usepackage{amssymb}
\usepackage{algorithm}
\usepackage{algorithmic}
\usepackage{booktabs}
\usepackage{multirow}
\usepackage{graphicx}
\usepackage{url}

%% Float placement. Twelve tables in a single narrow column will queue up and
%% drift far from the text that discusses them unless LaTeX is given room.
%% The defaults reserve most of a page for prose (textfraction 0.2) and allow
%% only two floats at the top, which is what pushed tables several pages past
%% their discussion. These are the standard relaxations, not forced placement:
%% the [H] specifier from the float package would pin every table exactly
%% where it appears in the source and leave large vertical gaps instead.
\renewcommand{\topfraction}{0.9}
\renewcommand{\bottomfraction}{0.8}
\renewcommand{\textfraction}{0.07}
\renewcommand{\floatpagefraction}{0.75}
\setcounter{topnumber}{3}
\setcounter{bottomnumber}{2}
\setcounter{totalnumber}{5}

\journalname{Empirical Software Engineering}

\begin{document}

\title{Cross-Corpus Evaluation of Domain-Aware Query Reformulation for
Regulatory Traceability}
\titlerunning{Cross-corpus evaluation of query reformulation}

\author{Sumanshu~Sohal \and Darshankumar~Prajapati}
\authorrunning{S. Sohal and D. Prajapati}

\institute{S. Sohal \at
              Department of Information Technology,
              University of the Cumberlands,
              Williamsburg, KY 40769, USA \\
              \email{sumanshu.95s@outlook.com}
           \and
           S. Sohal \at
              Independent Researcher, Washington, DC, USA
           \and
           D. Prajapati \at
              Independent Researcher, Old Bridge, NJ, USA \\
              \email{darshan3298@gmail.com}
}

% Left as the template supplies it. usrguid3 and template.tex both carry this
% literal string, and the template notes that the editor enters the real
% dates; it is not something the author fills in before submission.
\date{Received: date / Accepted: date}

\maketitle

\begin{abstract}
\textbf{Context.} Compliance mapping links regulatory requirements to the controls that address them. Public collections are scarce, so methods are usually measured on benchmarks their own authors build.

\textbf{Objective.} We test whether the measured benefit of domain-aware query reformulation generalizes from an author-constructed benchmark to three externally authored corpora.

\textbf{Method.} We score a frozen set of retrieval methods once per requirement over the full control corpus, on four corpora: one we built to isolate vocabulary mismatch, and three using externally authored framework or statutory texts linked through NIST crosswalks, sharing one 300-control target catalogue and so not independent of one another.

\textbf{Results.} Domain-aware query reformulation improves Top-1 accuracy by $+0.121$ on the corpus we built (95\% CI $[+0.034, +0.224]$, nominal $p = 0.039$). That estimate did not recur: the external corpora give $+0.028$, $0.000$ and $+0.032$, none distinguishable from zero. Each rests on very few requirements, flipping top-1 correctness for 9 of 58 on our corpus and 5, 4 and 5 externally, and each moves by up to $0.029$ under a backend-refitting choice unrelated to reformulation. We found no evidence that vocabulary distance moderates the effect within a corpus ($p = 0.31$). A preregistered gated hybrid shows the same spread, $+0.085$ on one corpus against $-0.088$ on another.

\textbf{Conclusions.} Measuring a method only on a benchmark its own authors built can overstate how far it transfers. At this effect size, the number of decisions behind an estimate, and its movement under unrelated fitting choices, are more informative than the estimate alone. Corpora, code and the preregistration are public.
\keywords{Regulatory compliance \and Requirements traceability \and Information retrieval \and Query reformulation \and Benchmark validity \and Cross-corpus evaluation}
\end{abstract}


%% ================================================================
\section{Introduction}
\label{sec:intro}
%% ================================================================

Regulatory compliance is a first-order concern for organizations operating in regulated industries.
Regulations, standards and frameworks, among them GDPR, HIPAA, PCI~DSS, NIST~CSF, ISO~27001, SOX and SOC~2, impose or organize requirements that must be demonstrably linked to implemented controls. That linking is known as \emph{compliance mapping}.
In practice, this mapping is overwhelmingly manual: compliance teams maintain spreadsheets linking regulatory clauses to control descriptions, a process that is labor-intensive, error-prone, and resistant to the pace of regulatory change.
The requirements engineering community has studied the difficulty of relating legal text to software artifacts for nearly two decades, from the extraction of rights and obligations from privacy regulations~\citep{Breaux2008} and the systematic treatment of legal texts as a source of requirements~\citep{OttoAnton2007}, to machine-learning support for tracing regulatory codes to product requirements~\citep{ClelandHuang2010}.
This body of work consistently reports that legal language is ambiguous, cross-referential, and written at a different level of abstraction than the artifacts it governs~\citep{Massey2010,Zeni2015}. These observations motivate the retrieval problem we address here.

From a software engineering perspective, compliance mapping is a form of \emph{traceability recovery}~\citep{Borg2014,Cleland-Huang2014}: establishing and maintaining links between documentation-like source artifacts (regulations, standards) and implementation-like target artifacts (controls, configurations, audit evidence).
However, compliance mapping introduces challenges absent from classical code--documentation traceability.
First, the \emph{open-world assumption}: the absence of a link is itself a meaningful signal, indicating a compliance gap that may require remediation.
Second, \emph{vocabulary mismatch}: regulatory text employs legal and policy-level language (e.g., ``implement appropriate technical measures including encryption of personal data'') that differs systematically from the implementation-level language of controls (e.g., ``AES-256 with FIPS 140-2 validated modules'').
Third, \emph{scope confusion}: lexically similar controls may address orthogonal regulatory concerns (e.g., encryption at rest vs.\ encryption in transit).
Fourth, \emph{regulatory drift}: both requirements and controls evolve independently over time~\citep{Gama2014}.

Existing information retrieval approaches to traceability recovery, whether based on vector space models~\citep{Antoniol2002}, Latent Semantic Indexing~\citep{MarcusMaletic2003}, or more recent embedding-based methods~\citep{Guo2017}, operate in a single-shot fashion: a query is issued, a ranked list is returned, and the top result is taken as the answer.
This paradigm is ill-suited to compliance mapping for two reasons.
First, when a single-shot retrieval returns a low-confidence result, no mechanism exists to refine the search or seek corroborating evidence.
Second, in high-stakes domains, the system should be able to \emph{abstain} rather than commit to a potentially incorrect mapping, a capability absent from standard retrieval pipelines.

In this paper, we introduce the \emph{Regulatory Alignment Agent} (RAA), a multi-step retrieval system that addresses these limitations through a deterministic, tool-augmented multi-stage workflow.
The design is inspired by the ReAct paradigm~\citep{Yao2023} of interleaving reasoning-like observations with tool actions, but, unlike an open-ended language-model ReAct loop, RAA executes a fixed conditional sequence of deterministic tools and emits structured (templated) trace steps: it retrieves candidate controls, evaluates confidence, reformulates queries using domain knowledge when the top-2 margin is ambiguous, corroborates candidates across related regulatory frameworks, and makes an accept-or-abstain decision recorded together with the scores and thresholds behind it. We use the term ``agent'' in this restricted, deterministic sense throughout.

\noindent\textbf{Research questions.}
The study is organized around three questions, stated here so the reader can check the results against them.

\begin{description}
\item[RQ1] Under one common ranking protocol, does the benefit of domain-aware query reformulation observed on a corpus we built also appear on corpora whose requirements, controls and links were authored outside this work?
\item[RQ2] \emph{Within} a corpus, does requirement-level vocabulary distance identify which requirements benefit from reformulation?
\item[RQ3] Does a gated combination of lexical and semantic retrieval, whose design was frozen and publicly timestamped before execution, recover what either family uniquely solves without giving up the other's successes?
\end{description}

RQ1 asks about generalization across corpora, not about corpus types as populations. We have four corpora, one of them ours, so the design can show whether an effect reappears elsewhere. It cannot attribute the difference to any single property of the corpus we built.

\noindent\textbf{Contributions.}
This paper makes four contributions.

First, a cross-corpus evaluation of regulatory traceability (Section~\ref{sec:dataset}) that puts one corpus we built alongside three whose links are NIST-authored crosswalks over externally authored framework or statutory text, all sharing a single 300-control target catalogue, and scores every method once per requirement over the full corpus so that all ranking comparisons share one estimand.

Second, the findings that answer RQ1 and RQ2 (Section~\ref{sec:results}). Reformulation improves Top-1 by $+0.121$ on our corpus and by $+0.028$, $0.000$ and $+0.032$ on the external ones, none of the three distinguishable from zero, and we found no evidence that a requirement-level vocabulary distance identifies who benefits within the external corpora. We report two properties of these estimates alongside them. Each turns on fewer than ten requirements whose top-1 correctness changes at all, and each moves by up to $0.029$ when the latent-semantic backend is refitted on a different population, a choice unrelated to reformulation. Since the mechanism of reformulation is a thesaurus written by the same authors as the corpus it succeeds on, this is a concrete case of a benchmark failing to predict how the intervention it exercises behaves elsewhere. We are explicit that a single corpus of our own construction cannot identify \emph{which} of its properties is responsible.

Third, a preregistered answer to RQ3 (Section~\ref{sec:hybrid}). The gated hybrid rules out a mean loss of five Top-1 points or more against semantic retrieval alone and shows no evidence of superiority, while producing a large gain on one corpus and a nearly equal loss on another. The design, its endpoint, its margin, its test order and its prohibitions were frozen, hashed and publicly registered before the analysis ran, which constrains outcome-dependent analysis and lets a result with no headline improvement be reported as evidence.

Fourth, a formalization of compliance mapping as a \emph{selective} traceability recovery problem (Section~\ref{sec:formulation}) that separates ranking quality from decision quality, together with the multi-step architecture used as the vehicle for the study (Section~\ref{sec:architecture}).

The system is the instrument here, not the claim. Once ranking is decoupled from the accept/abstain decision and compared at the requirement level, no added deterministic component significantly improves ranking on any external NIST-derived corpus, the agent holds no coverage or selective-accuracy advantage, and an open-world evaluation shows its abstention detects only a minority of constructed no-match cases.
We draw no compensating qualitative claim from this.
The traces record the steps behind each decision, which is a provenance property; whether they help an auditor is an empirical question this paper does not test.
All four corpora, the frozen analysis specification, the baselines and the implementation are public, and the preregistration is at \texttt{doi:10.17605/OSF.IO/NZXRV}.


%% ================================================================
\section{Related work}
\label{sec:related}
%% ================================================================

Our work sits at the intersection of several research threads: legal and regulatory requirements engineering, automated traceability recovery, query reformulation and rank fusion in information retrieval, selective prediction, and the methodological literature on what retrieval benchmarks license us to conclude.
We discuss each in turn and position our contribution against them.
The last thread matters most for this paper, because our central finding is about how a result behaves across corpora rather than about a system.

\subsection{Regulatory requirements and traceability recovery}
\label{sec:related-legal}

Treating legal texts as a source of software requirements has a long history in the requirements engineering (RE) community.
\citet{Breaux2008} developed a systematic methodology for extracting rights and obligations from privacy regulations such as HIPAA, and \citet{OttoAnton2007} surveyed the approaches available for addressing legal texts during requirements analysis, cataloguing the difficulties that regulation-derived requirements pose: ambiguity, cross-references, and evolving interpretations.
\citet{Massey2010} demonstrated these difficulties concretely by evaluating the security and privacy requirements of an electronic health record system for legal compliance, finding that even carefully specified requirements exhibited substantial compliance gaps.
Subsequent work sought tool support for this analysis: GaiusT~\citep{Zeni2015} semi-automatically annotates legal documents to extract actors, rights, and obligations; the N\`omos family of modelling languages~\citep{Ingolfo2011} represents law as models against which requirements compliance can be argued; and Legal-GRL~\citep{Ghanavati2014} adapts goal-oriented modelling to regulations.
\citet{Akhigbe2019} provide a systematic mapping of such modelling approaches for legal compliance.
A parallel thread in business process management addresses compliance of processes rather than requirements; \citet{Hashmi2018} survey this literature comprehensively.

Most directly related to our problem is the work of \citet{ClelandHuang2010}, who trained a probabilistic classifier to trace regulatory codes (HIPAA provisions) to product-level requirements, explicitly framing regulatory compliance as a traceability problem.
Our work shares this framing but differs in three respects: we map regulations to \emph{controls} across seven frameworks rather than to requirements of a single system; we adopt a multi-step retrieval architecture rather than a single-shot classifier; and we treat abstention as a first-class outcome.

More recently, NLP has been applied to compliance analysis at scale, notably by the requirements engineering group at the University of Luxembourg: automated extraction of semantic legal metadata~\citep{Sleimi2018}, question answering over compliance documents~\citep{Abualhaija2022}, completeness checking of privacy policies against GDPR~\citep{Amaral2022}, and compliance checking of data processing agreements~\citep{Cejas2023}.
Closest in task to ours, \citet{Etezadi2026} compare a similarity-based classifier with LLM prompting for tracing requirements to GDPR provisions, concluding that generic traceability techniques transfer poorly to legal artifacts. Our cross-corpus results reinforce that conclusion from the retrieval side.
Unlike these approaches, which rely on learned models, our agent is deliberately built from deterministic components, so every mapping decision is accompanied by a decision-provenance record listing the steps, scores and thresholds that produced it.

Compliance mapping is also a form of traceability recovery, and that literature is older still.
The application of information retrieval to software traceability has a two-decade history.
\citet{Antoniol2002} first demonstrated that probabilistic and vector space models could recover links between documentation and source code.
\citet{MarcusMaletic2003} subsequently showed that Latent Semantic Indexing (LSI) could capture latent relationships between artifacts even when surface terminology differed, a finding that motivated a generation of traceability tools.

Subsequent work has explored numerous extensions: incorporating structural information~\citep{McMillan2009}, leveraging developer feedback~\citep{Cuddeback2010}, applying web-mining techniques~\citep{Ali2013}, and tuning retrieval parameters~\citep{Oliveto2010}.
\citet{Borg2014} provide a comprehensive survey, noting that while IR-based approaches achieve reasonable recall, precision remains a persistent challenge, particularly for large artifact collections with high vocabulary overlap.

More recently, deep learning approaches have been applied to traceability.
\citet{Guo2017} used recurrent neural networks to learn distributed representations of artifacts.
Pre-trained language models such as BERT~\citep{Devlin2019} and domain-adapted variants (Legal-BERT~\citep{Chalkidis2020}, SecureBERT~\citep{Aghaei2022}) offer richer representations but require domain-specific fine-tuning.

Our work differs from these approaches in three ways: (1)~we address the \emph{compliance} variant of traceability where open-world non-matches (gaps) are first-class outcomes; (2)~we adopt a multi-step retrieval architecture rather than single-shot retrieval; and (3)~we evaluate the system's \emph{decision quality} (when to accept vs.\ abstain) alongside ranking quality.

\subsection{Retrieval techniques: reformulation, fusion, abstention and language models}
\label{sec:related-methods}

The agent is assembled from four families of technique, each with its own literature.

Query expansion and reformulation are well-established techniques for bridging vocabulary mismatch~\citep{Carpineto2012}.
Pseudo-relevance feedback~\citep{Rocchio1971} and thesaurus-based expansion~\citep{Voorhees1994} both aim to enrich the query with terms likely to appear in relevant documents, while neural rerankers such as BERT-based passage re-ranking~\citep{Nogueira2019} instead re-score candidates with a learned relevance model.

In specialized domains, curated thesauri have proven particularly effective.
Medical information retrieval benefits from UMLS-based expansion~\citep{Aronson2010}, and legal retrieval from legal terminology resources~\citep{Maxwell2017}.
To our knowledge, no prior work has applied systematic query reformulation to regulatory compliance mapping with a curated compliance-domain thesaurus.

Combining the outputs of multiple retrieval systems is a well-studied problem.
CombSUM, CombMNZ~\citep{Fox1994}, and Reciprocal Rank Fusion (RRF)~\citep{Cormack2009} are representative score- and rank-based fusion methods.
RRF is attractive for our setting because it requires no score normalization and is robust to heterogeneous scoring functions, which is the situation when combining BM25, TF-IDF and LSI outputs.

In settings where incorrect predictions carry high costs, selective prediction~\citep{ElYaniv2010,Geifman2017} allows a system to abstain on inputs where it is uncertain.
Coverage (the fraction of inputs on which the system makes a prediction) and selective accuracy (accuracy among accepted predictions) characterize the precision-coverage trade-off.
Selective prediction has been applied to classification~\citep{Geifman2019} and question answering~\citep{Kamath2020} but, to our knowledge, not to traceability recovery or compliance mapping.

The ReAct framework~\citep{Yao2023} demonstrated that interleaving reasoning traces with tool invocations enables language model agents to solve complex tasks through multi-step planning.
Subsequent work has applied agentic patterns to code generation~\citep{Yang2024}, scientific discovery~\citep{Bran2024}, and information retrieval~\citep{Shao2024}.
We adapt the ReAct pattern to compliance mapping, where the ``tools'' are retrieval backends, a domain thesaurus, and cross-framework knowledge. Not all of these are rule-based: latent semantic indexing and the pre-trained encoders are learned. What they share is that every component runs at fixed inference settings and exposes its inputs, transformations, scores and outputs, so each step can be recorded rather than inferred.

Language models have also been used as retrieval components directly rather than as planners.
Retrieval-augmented generation~\citep{Lewis2020} conditions generation on retrieved passages, and a line of work uses the model itself as a reranker: \citet{Sun2023} show that an instruction-following model can reorder a candidate list competitively with supervised cross-encoders, without task-specific training.
This matters for our setting because it changes what a reranking baseline is.
A model that reranks a candidate list inherits that list's ceiling exactly as a cross-encoder does, so the protocol distinction we draw in Section~\ref{sec:protocol} between end-to-end and candidate-constrained methods applies to LLM rerankers unchanged, and comparing a one-pass LLM against methods evaluated over repeated holdouts compares different estimands rather than different methods.

We do not evaluate a language model reranker here. This paper estimates the effect of one deterministic intervention and asks whether that estimate generalizes across corpora; it does not attempt to identify the strongest available regulatory retrieval system, and the methods it compares are those listed in Section~\ref{sec:protocol}. That is a scope boundary, and it carries a cost we state in Section~\ref{sec:threats}: our conclusions about what does and does not help are bounded by the method set we measured, and a prompted or instruction-tuned reranker could behave differently from every method here.

\subsection{Benchmark validity and the transfer of retrieval results}
\label{sec:related-benchmarks}

A recurring finding in information retrieval is that measured improvements need not survive a change of collection or of baseline.
\citet{Armstrong2009} examined a decade of ad-hoc retrieval results and found that reported gains largely failed to accumulate, because systems were repeatedly compared against weak baselines rather than against the strongest available.
\citet{Yang2019} reached a similar conclusion for early neural ranking models, showing that several reported gains shrank or disappeared against a properly tuned lexical baseline.
\citet{Voorhees2000} had earlier established a more precise point about the measurement instrument itself: assessor disagreement changes individual relevance judgments and absolute effectiveness scores, while the comparative ranking of systems remains substantially more stable.
The useful reading is therefore not that evaluation is unreliable, but that absolute figures belong to a collection and its assessors while comparisons between systems travel better.

The response in modern retrieval has been to evaluate across heterogeneous collections rather than to optimize one.
BEIR~\citep{Thakur2021} assembles retrieval tasks from different domains specifically to test zero-shot transfer, and reports that performance varies substantially across collections: BM25 remains a robust baseline, while reranking and late-interaction methods perform best on average at greater computational cost.
The lesson we take is not that one family of methods wins, but that a single collection is a poor guide to how a method will behave on the next one.

Adjacent fields have made a compatible point with a useful nuance.
\citet{Recht2019} rebuilt image classification test sets following the original collection procedure and found substantial drops in absolute accuracy, while the relative ordering of models and the gains between them largely transferred.
Sensitivity to how a dataset is constructed therefore shows up in levels more readily than in rankings, which is why we report per-corpus levels alongside contrasts rather than only the contrasts.
\citet{BowmanDahl2021}, in a position paper, argue that benchmarks in natural language understanding often fail to measure the capability they claim to, raising construct validity rather than reliability as the central concern.
Within software engineering, \citet{Sim2003} argued that shared benchmarks accelerate progress and help form research communities; we draw the further inference, which is ours rather than theirs, that this makes the construction of a benchmark a research contribution deserving the same scrutiny as a method.

This literature bears directly on our design.
Traceability research has often evaluated on a single collection, and the compliance-mapping variant has few public collections of any kind, which creates pressure to construct one.
A constructed collection is a legitimate diagnostic instrument, but its authors also choose its vocabulary, and that vocabulary may align with the mechanism of the intervention being tested, which would overestimate how far the intervention transfers.
We therefore evaluate the same frozen methods on one engineered corpus and three externally authored ones, whose source texts and mappings were not written by us, and report the contrast between them as the result rather than reporting the best number.

The three external corpora are independent of us but not of each other.
All three map to the same SP~800-53 control catalogue through NIST-authored crosswalks, so they are correlated evaluation corpora rather than independent samples of regulatory text, and we treat them as such throughout.
We are not aware of prior work in regulatory traceability that compares conclusions drawn from an engineered corpus against externally authored ones, though we make this claim from familiarity with the literature rather than from a structured search.
Section~\ref{sec:results} shows the contrast is large enough to reverse a conclusion drawn from the engineered corpus alone.


%% ================================================================
\section{Problem formulation}
\label{sec:formulation}
%% ================================================================


Let $\mathcal{R} = \{r_1, \ldots, r_n\}$ denote a set of regulatory requirements and $\mathcal{C} = \{c_1, \ldots, c_m\}$ a corpus of organizational controls.
The ground-truth mapping is a relation $\mathcal{G} \subseteq \mathcal{R} \times \mathcal{C}$, where $(r_i, c_j) \in \mathcal{G}$ indicates that control $c_j$ addresses (fully or partially) requirement $r_i$.
The mapping is many-to-many: a requirement may be addressed by multiple controls, and a control may satisfy multiple requirements.

A \emph{compliance mapping system} is a function $f: \mathcal{R} \rightarrow \mathcal{C}^k \times \{\texttt{accept}, \texttt{abstain}\}$ that, for each requirement $r$, produces a ranked list of $k$ candidate controls and a binary decision: \texttt{accept} the top-ranked mapping, or \texttt{abstain} and escalate for human review.

\label{sec:metrics}

We evaluate two orthogonal aspects of system quality:

\noindent\textbf{Ranking quality.}
For each requirement $r_i$ with ground-truth controls $\mathcal{G}_i = \{c : (r_i, c) \in \mathcal{G}\}$, we measure the quality of the ranked list using standard IR metrics: Top-1 accuracy (whether the highest-ranked control is correct), Mean Reciprocal Rank at $k$ (MRR@$k$), normalized Discounted Cumulative Gain at $k$ (nDCG@$k$), Mean Average Precision at $k$ (MAP@$k$), Recall@$k$, and Precision@$k$.

To ensure methodological continuity with the traceability literature, we additionally report \emph{legacy micro-precision and micro-recall}~\citep{MarcusMaletic2003}, computed by aggregating true positives, false positives, and false negatives across all queries:
\begin{equation}
\text{Precision}_\mu = \frac{\sum_i |\hat{\mathcal{G}}_i \cap \mathcal{G}_i|}{\sum_i |\hat{\mathcal{G}}_i|}, \quad
\text{Recall}_\mu = \frac{\sum_i |\hat{\mathcal{G}}_i \cap \mathcal{G}_i|}{\sum_i |\mathcal{G}_i|}
\end{equation}
where $\hat{\mathcal{G}}_i$ denotes the set of controls retrieved for requirement $r_i$ at a given cut point.

\noindent\textbf{Decision quality.}
For the selective prediction component, we report \emph{coverage} $\phi = |\{r : f(r).\text{decision} = \texttt{accept}\}| / |\mathcal{R}|$ and \emph{selective accuracy} $\text{SelAcc} = \text{Top-1 among accepted}$, following \citet{Geifman2017}.
The key design question is the precision--coverage trade-off: higher coverage is desirable (fewer escalations), but only if the accepted mappings are reliable.


%% ================================================================
\section{Approach: Regulatory Alignment Agent}
\label{sec:architecture}
%% ================================================================


The RAA processes each regulatory requirement $r$ through a fixed multi-stage sequence (Algorithm~\ref{alg:agent}).
The sequence is parameterized by a set of \emph{tools}: stateless, deterministic functions invoked in a fixed order and gated by scale-invariant confidence tests rather than by open-ended planning.
Each invocation produces an \texttt{AgentStep} record $(s, a, \text{input}, o)$ consisting of a templated status string $s$, an action $a$, its input, and the resulting observation $o$.
The complete sequence of steps forms an \texttt{AgentTrace}, the decision-provenance record for that mapping. Whether such a record is useful to an auditor is untested here.

\begin{algorithm}[t]
\caption{RAA execution loop}
\label{alg:agent}
\begin{algorithmic}[1]
\REQUIRE Requirement $r$, backends $\mathcal{B}$, thesaurus $\mathcal{T}$, thresholds $\theta_c, \theta_g$
\ENSURE Decision $d \in \{\texttt{accept}, \texttt{abstain}\}$, ranked list, trace
%
\STATE $\mathbf{s}_{\text{best}} \gets \textsc{RRF}(\{\textsc{Retrieve}(r.\text{text}, b) : b \in \mathcal{B}\})$ \COMMENT{Multi-backend fusion}
\STATE $\text{conf}, \text{gap} \gets \textsc{TopConfGap}(\mathbf{s}_{\text{best}})$
%
\IF{$\text{gap} / \text{conf} < \tau_{\text{rel}}$}  \COMMENT{Ambiguous top-2 margin (scale-invariant)}
    \STATE $q' \gets \textsc{Reformulate}(r.\text{text}, \mathcal{T})$ \COMMENT{Query expansion}
    \STATE $\mathbf{s}_{\text{ref}} \gets \textsc{RRF}(\{\textsc{Retrieve}(q', b) : b \in \mathcal{B}\})$
    \IF{$\textsc{TopConf}(\mathbf{s}_{\text{ref}}) > \text{conf}$} \STATE $\mathbf{s}_{\text{best}} \gets \mathbf{s}_{\text{ref}}$ \ENDIF
    \IF{decomposition enabled \AND $r$ is compound}
        \STATE $\mathbf{s}_{\text{best}} \gets \textsc{Decompose\&Aggregate}(r, \mathcal{B})$ \COMMENT{Optional; own ablation flag}
    \ENDIF
\ENDIF
%
\STATE $c^* \gets \arg\max_j \mathbf{s}_{\text{best}}[j]$; \quad $\theta'_c \gets \theta_c$, $\theta'_g \gets \theta_g$
\STATE $\rho \gets \textsc{CrossRef}(r.\text{framework}, c^*.\text{family})$ \COMMENT{Corroboration (decision only)}
\STATE $\theta'_c \gets \theta'_c \cdot (1 - 0.15\rho)$ \COMMENT{Relax threshold if corroborated}
\IF{verification enabled \AND \NOT $\textsc{ReverseRankOK}(r, c^*)$}
    \STATE $\theta'_c \gets \theta'_c \cdot 1.1$, \quad $\theta'_g \gets \theta'_g \cdot 1.1$ \COMMENT{Tighten if inconsistent}
\ENDIF
\STATE $\text{conf}, \text{gap} \gets \textsc{TopConfGap}(\mathbf{s}_{\text{best}})$ \COMMENT{Ranking never reordered by CrossRef/Verify}
%
\IF{$\text{conf} \geq \theta'_c$ \AND $\text{gap} \geq \theta'_g$}
    \RETURN $\texttt{accept}$, ranked list, trace
\ELSE
    \RETURN $\texttt{abstain}$, ranked list, trace
\ENDIF
%\end{algorithmic}
\end{algorithm}


The workflow draws on tools that fall into two groups: those that shape the \emph{ranking} of candidate controls (fusion, reformulation, optional decomposition) and those that shape only the \emph{decision} to accept or abstain (corroboration, verification, calibrated thresholding). Keeping the two groups separate is deliberate: it lets us attribute ranking changes and decision changes to distinct components, and it prevents a decision heuristic from silently reordering results.

The first ranking tool is \emph{multi-backend retrieval and fusion}.
The agent maintains three retrieval backends (TF-IDF with cosine similarity, Okapi BM25~\citep{Robertson2009}, and Latent Semantic Indexing~\citep{MarcusMaletic2003}), each producing a score vector $\mathbf{s}_b \in \mathbb{R}^m$ over the control corpus.
Their rankings are combined via Reciprocal Rank Fusion~\citep{Cormack2009},
\begin{equation}
\text{RRF}(c_j) = \sum_{b \in \mathcal{B}} \frac{1}{k + \text{rank}_b(c_j)},
\end{equation}
with the standard smoothing constant $k = 60$.
RRF requires no score normalization and is robust to the heterogeneous scoring functions of the three backends.

The second ranking tool, \emph{domain-aware query reformulation}, is invoked when the top-2 result is \emph{ambiguous}, defined by a scale-invariant criterion: the relative margin $\text{gap}/\text{conf}$ between the top two fused scores falls below a threshold $\tau_{\text{rel}}$ (0.10 in our experiments). We use a relative rather than absolute margin because backend score scales differ by orders of magnitude, raw BM25 scores and fused RRF scores being on quite different footings, so a single absolute confidence threshold cannot be shared across configurations. Reformulation addresses vocabulary mismatch through two deterministic mechanisms: a set of 32 regular-expression patterns that map regulatory phrases to implementation concepts (for instance, ``boundary protection'' expands to firewall, network perimeter, DMZ, segmentation, and IPS), and a curated thesaurus of 20 concept families that supplies synonym and hypernym expansions.
The full thesaurus is reproduced in \ref{app:thesaurus}.
The reformulated query is the concatenation of the original query and the deduplicated expansion terms; no neural model is involved, so each expansion records the terms that produced it and the pattern or thesaurus entry they came from.
An optional third ranking tool, \emph{query decomposition}, splits a compound requirement into sub-clauses, retrieves for each, and averages the results. Because decomposition is logically distinct from reformulation, we expose it under its own ablation flag rather than bundling the two.

Turning to the decision tools, \emph{cross-framework corroboration} exploits the fact that regulatory frameworks overlap: NIST and ISO share access-control concepts, and HIPAA and SOC~2 overlap on audit requirements.
Using a hand-coded framework relationship graph and a taxonomy grouping controls by functional concern, the tool checks whether the top candidate's domain family spans frameworks related to the requirement's framework, yielding a corroboration score $\rho \in [0, 1]$.
Importantly for leakage, this signal is derived only from the static family-to-framework taxonomy, never from the ground-truth mapping; when $\rho > 0$ the agent relaxes its acceptance threshold by up to 15\%, admitting borderline mappings that have structural support.
\emph{Bidirectional verification} then queries the requirement corpus with the top control's own text and locates the original requirement in that reverse ranking, in the spirit of cycle-consistency checks in unsupervised machine translation~\citep{Lample2018}. If the requirement is not recovered within the top decile of the reverse ranking, the mapping is deemed inconsistent and the acceptance thresholds are tightened by 10\%, nudging the borderline case toward abstention. Verification never reorders the candidate list; it acts purely on the decision.
Finally, the \emph{calibrated selective decision} governs the accept/abstain outcome through a confidence threshold $\theta_c$ and a score-gap threshold $\theta_g$, both calibrated on a held-out split to target a specified coverage level; the calibration procedure is detailed in Section~\ref{sec:protocol}.


%% ================================================================
\section{Evaluation corpora}
\label{sec:dataset}
%% ================================================================

We evaluate on four corpora: one we constructed, and three assembled entirely from artifacts published by NIST, in which both the texts and the ground-truth mappings were authored independently of this work.
The engineered corpus tells us \emph{why} a component helps, because we control the degree of vocabulary mismatch in it.
The three external corpora tell us whether that help survives contact with externally authored framework and statutory text, and, because there are three rather than one, whether a conclusion drawn from any single one of them would have held on the others.

Table~\ref{tab:corpora} summarizes them. All three external corpora map into the same 300-control corpus drawn from SP~800-53, so differences between them are differences in the regulatory side of the mapping rather than in the target catalogue.

\begin{table}[htbp]
\caption{The four evaluation corpora. Recall@20 is the share of requirements for which the dual-encoder places a gold control in its top 20, and bounds any method that reranks that shortlist.}
\label{tab:corpora}
\centering
\setlength{\tabcolsep}{3pt}
\small
\resizebox{\ifdim\width>\columnwidth\columnwidth\else\width\fi}{!}{%
\begin{tabular}{@{}llrrrc@{}}
\toprule
Corpus & Ground truth & Reqs & Ctrls & Links & R@20 \\
\midrule
Diagnostic & authored here & 58 & 110 & 86 & .983 \\
CSF~1.1 & NIST crosswalk & 106 & 300 & 495 & .906 \\
HIPAA & NIST OLIR & 68 & 300 & 274 & .882 \\
Privacy Fw. & NIST crosswalk & 94 & 300 & 456 & .840 \\
\bottomrule
\end{tabular}}
\end{table}

Two properties of this design need stating before the results.

The three external corpora are independent of us but not of each other. All map into one catalogue through NIST-authored crosswalks, so they are correlated evaluation corpora rather than independent samples of regulatory text, and a result that holds across all three is weaker evidence than three unrelated corpora would provide. We treat this as a limitation throughout and return to it in Section~\ref{sec:threats}.

NIST describes these mappings as concept relationship mappings rather than assertions of implementation, and marks the HIPAA OLIR ``Comprehensive: No''~\citep{NISTIR8477}. They are therefore a silver standard. A predicted link that the crosswalk does not list is not automatically wrong, and we say ``addresses'' rather than ``implements'' throughout.

\subsection{Diagnostic benchmark}
\label{sec:dataset-synth}

The diagnostic benchmark contains 58 requirements drawn from seven frameworks: GDPR (8), NIST CSF and SP~800-53 (12), HIPAA (10), PCI~DSS (10), ISO~27001 (8), SOX ITGC (4), and SOC~2 (6).
Each requirement is a single-clause statement we wrote to summarize an obligation in the source framework.
For ISO~27001, PCI~DSS and SOC~2 this is a paraphrase rather than a quotation, because those documents are not openly published, and we redistribute none of them.
The requirement vocabulary in this corpus is therefore ours, which is one reason it cannot settle whether a thesaurus we also wrote transfers elsewhere.
The 110-control corpus is stratified by design.
Fifty-nine controls are labelled \texttt{perfect}: near-paraphrases of the corresponding requirement that represent the easy retrieval case.
Twenty-seven are labelled \texttt{good}: correct mappings expressed in different terminology.
These are construction labels recorded when the corpus was written, not measurements of lexical overlap, and the two do not coincide.
The \texttt{good} stratum averages 0.043 IDF-weighted content-word overlap with its requirement against 0.250 for \texttt{perfect}, so the labels do track vocabulary distance, but 14 of the 59 \texttt{perfect} links also share no content words with their requirement.
For example, the GDPR Art.~32 encryption requirement is matched by ``Cryptographic safeguards applied to individually identifiable records stored in relational databases,'' a control sharing no content words with the requirement.
The remaining 24 controls are negatives: 20 \emph{hard negatives} spanning scope confounders (in-transit encryption as a distractor for at-rest encryption requirements), cross-family distractors, and near-miss negatives (development-only logging), plus four unrelated controls.
These counts are the \texttt{match\_type} labels in the released \texttt{diag\_controls.csv}, where the two positive strata appear as \texttt{perfect} and \texttt{good}.
The ground truth contains 86 positive links and is multi-label: 15 requirements map to two or more controls.

Because the authors constructed both the control corpus and the domain thesaurus used by the reformulation tool, results on this benchmark alone could overstate the value of reformulation: the mismatched controls and the thesaurus could encode the same design intuitions.
The three externally authored corpora below are what breaks that circle.

\subsection{The three externally authored corpora}
\label{sec:dataset-real}

The remaining three corpora use externally authored framework or statutory texts, linked through NIST-authored crosswalks, and all three map into one 300-control corpus drawn from SP~800-53r5. Two of the three requirement sets are NIST's own; the HIPAA Security Rule is not, being statutory text issued by the U.S.\ Department of Health and Human Services, and what NIST supplies there is the crosswalk. They are externally authored but not independent of one another, which Section~\ref{sec:threats} treats as the study's most important limitation.

\noindent\textbf{NIST CSF 1.1 to SP 800-53r5.}

The second corpus is derived from three artifacts published by NIST.
The requirement side consists of the 108 subcategories of the NIST Cybersecurity Framework version~1.1, each a single outcome-oriented statement (e.g., ``ID.AM-3: Organizational communication and data flows are mapped'').
The control side consists of the 300 non-withdrawn base controls of NIST SP~800-53 Revision~5, with each control's title and statement text extracted from the official OSCAL catalog and organization-defined parameters resolved to their assignment labels.
Ground truth is NIST's official mapping between CSF subcategories and SP~800-53r5 controls, published alongside SP~800-53r5.
Two subcategories whose informative reference is ``all controls'' are excluded as unusable for retrieval evaluation, leaving 106 requirements and 495 links (mean 4.7 mapped controls per requirement, median 4, maximum 16).

This corpus is a substantially harder and more faithful test than the diagnostic benchmark in three respects.
The candidate pool is nearly three times larger (300 vs.\ 110 controls); the mapping is genuinely many-to-many with dense multi-label structure; and, critically, neither the texts nor the links were authored by us, so no component of the agent can be inadvertently tuned to the ground truth.
It also differs in character: CSF subcategories and SP~800-53 controls are written in the same institutional style by the same body, where the diagnostic benchmark deliberately puts policy language against implementation language. We describe that difference qualitatively and do not quantify it. Our gap measure is fitted per corpus and gives no scale on which one corpus can be called narrower-gapped than another.
An earlier version of this work had only this corpus and the diagnostic one, and concluded from the contrast between them that gap severity determines how much reformulation contributes. Sections~\ref{sec:results-multi} and \ref{sec:results-gap} test that conclusion against two further corpora and do not support it.
One caveat bounds the interpretation of this corpus: NIST characterizes its crosswalks as \emph{concept relationship} mappings between documentary standards, meaning that a control's outcomes support a CSF subcategory in whole or in part, rather than as evidence that any organization has implemented a control satisfying a requirement~\citep{NISTIR8477}.
Accordingly, we treat this benchmark as a test of reference-link recovery between framework artifacts, not of implemented-compliance verification, and we scope our claims to that task.
The generated CSV files are in Online Resource~1 and the extraction scripts are in the code repository, so the corpus can be rebuilt from the NIST sources.

\noindent\textbf{HIPAA Security Rule to SP 800-53.}
\label{sec:dataset-hipaa}

The third corpus maps requirements of the HIPAA Security Rule to SP~800-53 controls through NIST's OLIR crosswalk published with SP~800-66 Revision~2. It contains 68 requirements and 274 links.

It is included because its requirement side is statutory text rather than framework prose, written by a different institution for a different purpose than the control catalogue it maps into. It is the only corpus here in which the two sides of the mapping come from different institutions, which is the qualitative condition under which a regulation-to-implementation thesaurus should be most useful. We say ``qualitative'' deliberately: our gap measure is corpus-relative and gives no common scale on which to rank corpora by vocabulary distance, so we make no quantitative claim that this corpus has a wider gap than the others.

The builder fetches and hashes each of the six raw responses from NIST's Cybersecurity and Privacy Reference Tool and stores them alongside the generated files, so the corpus can be rebuilt without querying a live service. Every control reference resolves.

One version caveat. The OLIR targets SP~800-53 Revision~5.1.1 while our control corpus is drawn from 5.2.0. The two differ by a single control, SA-24, and all 1{,}241 links in the source crosswalk resolve in both, so we use 5.2.0 for consistency with the other corpora and note the substitution.

\noindent\textbf{NIST Privacy Framework to SP 800-53.}
\label{sec:dataset-pf}

The fourth corpus maps subcategories of the NIST Privacy Framework to the same control corpus, 94 requirements and 456 links, through the crosswalk NIST publishes alongside SP~800-53.

Its role is to add a second NIST-authored corpus with different subject matter, so that a null on CSF is not read as a property of one crosswalk. We deliberately do not describe either corpus as narrow-gap: that would be a between-corpus comparison our measure cannot support. The Privacy Framework is NIST-authored like CSF, but it is a different framework with a different subject matter, and its subcategory identifiers and texts have no overlap with the CSF corpus. Six subcategories that carry no link and five blanket references are dropped by the builder, control enhancements are collapsed to their base control, and no control reference is left unresolved.


%% ================================================================
\section{Experimental protocol}
\label{sec:protocol}
%% ================================================================

\subsection{Two measurements, and which one carries the ranking claims}
\label{sec:protocol-estimands}

The study reports two kinds of number, and they are produced differently.

\emph{Ranking} results are measured once per requirement over the full control corpus. Every requirement is scored exactly once, no split is involved, and the reported figure is the mean over all requirements of the corpus. This is the primary protocol, and every cross-corpus ranking comparison in the paper uses it.

\emph{Decision} results (coverage, selective accuracy, open-world gap detection) need an acceptance threshold, and a threshold has to be chosen on data the system is not then judged on. These come from repeated stratified holdouts: for each of 30 random seeds, every stratum contributes proportionally to three disjoint splits (training, calibration, test) with holdout ratio 0.20 and calibration ratio 0.15, stratified by regulatory framework family and, on the CSF corpus, by CSF function. Requirements are tested an uneven number of times across seeds, so each is collapsed to its mean over the splits it appeared in before any comparison.

The two are different estimands and we never mix them in one comparison. Tables carry the protocol in their captions.

We also fix where the latent-semantic backend is fitted, because that turns out to matter more than the protocol does. Under the primary protocol, LSI is fitted on the 300 control documents alone and requirement text never enters the fit. The repeated-holdout runs fit it on the controls plus the training and calibration requirements, which excludes the test requirements but is still a larger and different fitting population. Section~\ref{sec:results-lsi} separates the two effects and reports what each contributes.

\subsection{Variants, baselines and inference}

To isolate each component's marginal contribution, we evaluate six incrementally constructed variants: \textbf{Single} (one-shot BM25 retrieval with threshold-based accept/abstain), \textbf{+Multi} (adding Reciprocal Rank Fusion across TF-IDF, BM25, and LSI), \textbf{+Reform} (adding domain-aware query reformulation on ambiguous queries), \textbf{+Decomp} (adding query decomposition, disclosed as its own step rather than bundled with reformulation), \textbf{+CrossRef} (adding cross-framework corroboration), and the \textbf{Full Agent} (adding bidirectional verification).
We further compare against three lexical baselines (TF-IDF, BM25, LSI) evaluated independently, a neural dual-encoder (\texttt{all-MiniLM-L6-v2}~\citep{Reimers2019}) using cosine similarity over sentence embeddings, and a cross-encoder reranker (\texttt{ms-marco-MiniLM-L-6-v2}) applied to the top-20 dual-encoder candidates.
Because ranking-shaping and decision-shaping tools are separated in the architecture (Section~\ref{sec:architecture}), the +CrossRef and Full Agent rows differ from +Decomp only in their decision profile, never in the ranking they produce.

Each agent variant is calibrated \emph{end-to-end}: its base acceptance thresholds are chosen by running the complete pipeline on the calibration split, including reformulation and the corroboration and verification threshold adjustments, so that the reported operating point reflects the target coverage (0.80) rather than the coverage of an intermediate, pre-adjustment ranking. Repeated-holdout decision analyses use 30 seeds; the primary ranking results are computed in a single full-population pass.

\label{sec:protocol-stats}
Three further points about the statistics, all prompted by weaknesses we identified in earlier versions of this work.

First, all ranking metrics (Top-1, MRR, nDCG, recall, precision) are computed from the full candidate ranking for every query, independent of the accept/abstain decision; coverage, selective accuracy, and open-world gap detection are reported as separate decision metrics. Computing ranking metrics only over accepted queries conflates ranking effectiveness with coverage and can report a pure coverage change as a spurious ranking gain.

Second, we do \emph{not} treat the 30 seeds as independent replicates: the stratified splits overlap heavily, so an across-seed $t$-test is pseudoreplicated and overstates significance. All significance claims are made at the requirement level. For a given metric and pair of methods, we compare the two over the shared requirements with a paired sign-flip randomization test for the two-sided $p$-value and a paired bootstrap over requirements for the 95\% confidence interval of the mean difference. Seed-averaged means and their intervals still appear in the holdout tables as description, but no inference rests on them. Significance is assessed at $\alpha = 0.05$.

Third, we report wins, losses and ties with every paired effect. Top-1 is binary per requirement, so a paired difference is $+1$, $-1$ or $0$, and the mean difference is entirely determined by the requirements on which the two methods disagree. On corpora of this size that set is often smaller than ten. Reporting it is the only way for a reader to see how much of an estimate rests on how few decisions, and it is a fact about the data that no interval conveys on its own.

Equivalence against a pre-specified $\pm 0.05$ margin is reported as a secondary result using two one-sided $t$ tests. We flag any equivalence decision whose $p$ falls within $0.01$ of $\alpha$, because with so few discordant requirements such a decision can turn on a single one. We did not implement an exact matched-pair equivalence procedure; the exact sign tests reported alongside address superiority only, and say nothing about whether an effect lies inside a margin.

The open-world evaluation of Section~\ref{sec:results-openworld} additionally holds out a random 50\% of controls per seed to manufacture no-match queries.

\medskip\noindent\textbf{Use of AI tools in this study.} No component of the system under evaluation is a \emph{generative} language model, and no result reported here was produced by one. Two baselines use pre-trained transformer models at pinned revisions: a dual-encoder that produces fixed embeddings, and a cross-encoder that directly scores requirement--control pairs. Both are language models in the representational sense, and neither generates text. The rest of the pipeline is lexical backends, a hand-written thesaurus and regular-expression patterns, all at fixed inference settings. Generative AI assistants were used during the conduct of the work in two roles, both of which we state so that a reader can judge them. They were used to review and refactor the evaluation harness and the analysis scripts, including the provenance tooling described above, and to review and edit the prose of this manuscript. Every statistical procedure, number, table and citation was independently verified by the authors against the committed records, and the authors take full responsibility for the content.


%% ================================================================
\section{Results}
\label{sec:results}
%% ================================================================

\subsection{Cross-corpus comparison, and the per-corpus ablations}
\label{sec:results-unified}

Table~\ref{tab:unified} is the study's single like-for-like comparison. Every method is scored once per requirement over the full control corpus, with LSI fitted on controls only, so the four columns and the seven rows are all on the same footing.

\begin{table}[htbp]
\caption{Top-1 accuracy for every method on all four corpora under the primary protocol: one scoring pass per requirement over the full control corpus, LSI fitted on controls only. The lower block reranks the dual-encoder's top 20 rather than the corpus, so it is bounded by that shortlist; its ceiling is given beneath it.}
\label{tab:unified}
\centering
\setlength{\tabcolsep}{3pt}
\small
\resizebox{\ifdim\width>\columnwidth\columnwidth\else\width\fi}{!}{%
\begin{tabular}{@{}lrrrr@{}}
\toprule
Method & Diagnostic & CSF~1.1 & HIPAA & Privacy Fw. \\
\midrule
TF-IDF & .500 & .377 & .221 & .372 \\
BM25 & .466 & .358 & .221 & .372 \\
LSI & .448 & .368 & .221 & .362 \\
Lexical RRF & .466 & .377 & .309 & .340 \\
Full RAA & .586 & .406 & .309 & .372 \\
Dual-encoder & .621 & .406 & .456 & .372 \\
\midrule
\multicolumn{5}{@{}l}{\emph{Candidate-constrained: reranks the dual-encoder top-20}} \\
Cross-encoder & .586 & .547 & .412 & .457 \\
\quad Recall@20 ceiling & .983 & .906 & .882 & .840 \\
\bottomrule
\end{tabular}}
\end{table}

Three things follow, and they replace claims that were previously assembled across separate tables.

The deterministic agent never out-ranks the dual-encoder. Paired at the requirement level, full RAA minus dual-encoder is $-0.034$ on the diagnostic corpus ($p = 0.82$), $0.000$ on CSF, $-0.147$ on HIPAA ($p = 0.064$) and $0.000$ on the Privacy Framework. Two of those are exact aggregate ties, and they are not agreement: on CSF the two methods disagree on 22 requirements, 11 each way, and on the Privacy Framework on 28, 14 each way. Equal means conceal complementary errors, which is the observation that motivated the gated hybrid of Section~\ref{sec:hybrid}.

The two-stage cross-encoder pipeline is not uniformly better either, which a single-corpus reading would have missed. It is well ahead on CSF ($.547$ against $.406$) and the Privacy Framework ($.457$ against $.372$), and behind the dual-encoder alone on the diagnostic corpus ($.586$ against $.621$) and on HIPAA ($.412$ against $.456$). As complete systems these are fair comparisons and we report them as such. As a comparison between reranking methods the lower block is bounded by its shortlist, which is why the ceiling sits beneath it.

And the lexical methods cluster tightly on three of the four corpora. On the Privacy Framework TF-IDF, BM25 and the dual-encoder all score exactly $.372$, which is the tie discussed in Section~\ref{sec:hybrid}.

\noindent\textbf{Diagnostic benchmark.}
\label{sec:results-synth}
The remaining tables in this subsection are the repeated-holdout ablations. They carry the coverage and selective-accuracy columns, which need a calibrated threshold and therefore cannot come from the one-pass protocol. Their Top-1 columns are descriptive.

Table~\ref{tab:baselines} presents the lexical and neural baseline results on the diagnostic benchmark.

\begin{table}[htbp]
\caption{Baseline comparison on the diagnostic benchmark. Repeated-holdout protocol, mean $\pm$ 95\% CI over 30 seeds, 110 controls, 86 positive links. Descriptive; the cross-corpus ranking claims of Section~\ref{sec:results-multi} use the one-pass protocol instead.}
\label{tab:baselines}
\centering
\setlength{\tabcolsep}{3pt}
\small
\resizebox{\ifdim\width>\columnwidth\columnwidth\else\width\fi}{!}{%
\begin{tabular}{@{}lcccccc@{}}
\toprule
Method & Top-1 & MRR@5 & nDCG@5 & R@5 & Cov. & SelAcc \\
\midrule
TF-IDF        & .572 $\pm$ .041 & .662 $\pm$ .037 & .614 $\pm$ .036 & .666 $\pm$ .041 & .803 & .643 \\
BM25          & .533 $\pm$ .037 & .623 $\pm$ .035 & .578 $\pm$ .037 & .625 $\pm$ .042 & .739 & .656 \\
LSI~\citep{MarcusMaletic2003} & .511 $\pm$ .041 & .640 $\pm$ .034 & .601 $\pm$ .035 & .676 $\pm$ .040 & .922 & .542 \\
\midrule
Dual-enc.~\citep{Reimers2019} & .672 $\pm$ .044 & .762 $\pm$ .033 & .729 $\pm$ .034 & .800 $\pm$ .038 & .758 & .783 \\
Cross-enc.    & .647 $\pm$ .036 & .733 $\pm$ .031 & .692 $\pm$ .033 & .759 $\pm$ .034 & .719 & .779 \\
\bottomrule
\end{tabular}}
\end{table}

All ranking metrics are computed from the full candidate ranking, independent of the accept/abstain decision, so that ranking quality and decision quality are never conflated; coverage and selective accuracy are reported separately as decision metrics.

Three observations warrant discussion.
First, lexical baselines plateau near 0.51--0.57~Top-1, with TF-IDF leading.
This is consistent with the lexical structure of the corpus: 41\% of positive links (35/86) share no content words with the corresponding requirement, creating a retrieval ceiling for term-matching methods.
That share exceeds the 27 links labelled \texttt{good}, because 14 \texttt{perfect} links are also lexically disjoint.
Measured mismatch is therefore wider than the construction labels alone would suggest, and 35/86 rather than 27/86 is the figure that bears on term-matching methods.
Second, LSI achieves competitive Recall@5 (0.676) at the highest coverage (0.922) but the lowest selective accuracy (0.542): its latent space aids recall but yields poorly calibrated confidence scores, a liability for a system that must decide when to escalate.
Third, the neural encoders are the strongest rankers on this benchmark: the dual-encoder reaches 0.672~Top-1 and 0.762~MRR@5, and even the general-purpose cross-encoder (0.647~Top-1) exceeds every lexical baseline. Pre-trained sentence embeddings clearly capture the cross-vocabulary similarity that term matching misses.

Table~\ref{tab:ablation} presents the incremental agent ablation on the same benchmark.

\begin{table}[htbp]
\caption{Agent ablation on the diagnostic benchmark. Repeated-holdout protocol, mean $\pm$ 95\% CI over 30 seeds; $\Delta$ is the marginal Top-1 gain over the previous row. Ranking metrics are decision-independent. The coverage and selective-accuracy columns are the reason this table uses holdouts: they need a calibrated threshold. The Top-1 column is descriptive here; Table~\ref{tab:multi} gives the primary figure.}
\label{tab:ablation}
\centering
\setlength{\tabcolsep}{3pt}
\small
\resizebox{\ifdim\width>\columnwidth\columnwidth\else\width\fi}{!}{%
\begin{tabular}{@{}lccccccc@{}}
\toprule
Variant & Top-1 & $\Delta$ & MRR@5 & R@5 & Cov. & SelAcc & Steps \\
\midrule
Single        & .533 $\pm$ .037 & ---    & .623 $\pm$ .035 & .625 $\pm$ .042 & .739 & .656 & 2.0 \\
+Multi        & .550 $\pm$ .044 & +.017  & .655 $\pm$ .036 & .683 $\pm$ .039 & .789 & .644 & 3.0 \\
+Reform       & .644 $\pm$ .035 & +.094  & .737 $\pm$ .026 & .738 $\pm$ .033 & .889 & .691 & 4.9 \\
+Decomp       & .644 $\pm$ .035 & .000   & .737 $\pm$ .026 & .738 $\pm$ .033 & .889 & .691 & 5.0 \\
+CrossRef     & .644 $\pm$ .035 & .000   & .737 $\pm$ .026 & .738 $\pm$ .033 & .844 & .639 & 6.0 \\
Full Agent    & .644 $\pm$ .035 & .000   & .737 $\pm$ .026 & .738 $\pm$ .033 & .833 & .648 & 7.0 \\
\bottomrule
\end{tabular}}
\end{table}

Statistical claims here use the requirement-level paired tests of Section~\ref{sec:protocol-stats}, which avoid the seed-level pseudoreplication inherent to overlapping splits.
The ablation yields four findings.
First, query reformulation is the only component that improves ranking. Under the primary one-pass protocol, +Reform improves Top-1 over +Multi by $+0.121$ (95\% bootstrap CI $[+0.034, +0.224]$; sign-flip randomization $p = 0.039$), and it does so by flipping top-1 correctness on 9 of the 58 requirements, winning on 8 and losing on 1. It reorders the top-ranked control on 15 and changes the ranking at some position on 21, so most of what it does is invisible in the mean. The seed-averaged holdout figures in the table are larger ($0.550$ to $0.644$) because that protocol also fits LSI on the training and calibration requirements; Section~\ref{sec:results-lsi} separates the two.
On this benchmark, where 41\% of the positive links (35/86) share no content words with their requirement, vocabulary mismatch is the main bottleneck and domain-aware reformulation is an effective mitigation for the requirements it reaches.
Second, multi-backend fusion does \emph{not} yield a significant ranking gain at the requirement level (Top-1 $+0.010$, $p = 0.88$ on the holdout data behind this table; the seed-averaged means nudge up from 0.533 to 0.550, but the effect does not survive a requirement-level test).
Third, the components beyond reformulation do not improve ranking at all, and we separate them rather than switching them on together. Adding decomposition alone changes the ranking for one requirement across all four corpora, HIPAA requirement~58, where it selects a different control; both the old and the new choice are wrong, so no reported figure moves. Adding corroboration and verification on top of decomposition changes no ranking on any requirement of any corpus, which is what the architecture intends and what we can now state as a measurement rather than an assertion. An earlier version of this work compared the full agent directly against the reformulation-only arm and read the single changed requirement as evidence that a decision-only tool had reordered. That contrast switches on three components at once, and the responsible one is decomposition. Their effect is confined to the decision profile: relative to +Reform (coverage 0.889), corroboration and verification trade coverage (down to 0.833) for calibration, the full agent abstaining more often on mappings that lack corroboration or fail the bidirectional check.
Fourth, the deterministic agent does not out-rank the neural baselines: the dual-encoder reaches 0.672 Top-1 and the cross-encoder 0.647, both nominally above the agent's 0.644, and at the requirement level the agent is statistically indistinguishable from each (agent vs.\ dual-encoder $p \approx 1.0$; agent vs.\ cross-encoder $p = 0.68$). Nor does the agent hold a decision-quality advantage: its coverage (0.833) is below LSI's (0.922) and its selective accuracy (0.648) below the dual-encoder's (0.783). The agent therefore shows no measured edge on this benchmark. We return to what does and does not follow from that in Section~\ref{sec:discussion}.

For methodological continuity with the traceability literature, Table~\ref{tab:legacy} additionally reports micro-precision and micro-recall at cut points 1 and 5, following the aggregation protocol of \citet{MarcusMaletic2003}.

\begin{table}[htbp]
\caption{Legacy micro-precision/recall at cut points on the diagnostic benchmark. Repeated-holdout protocol, 30 seeds; descriptive.}
\label{tab:legacy}
\centering
\setlength{\tabcolsep}{3pt}
\small
\resizebox{\ifdim\width>\columnwidth\columnwidth\else\width\fi}{!}{%
\begin{tabular}{@{}lcccc@{}}
\toprule
Method & $\text{P}_\mu$@1 & $\text{R}_\mu$@1 & $\text{P}_\mu$@5 & $\text{R}_\mu$@5 \\
\midrule
TF-IDF        & .572 $\pm$ .041 & .391 $\pm$ .035 & .174 $\pm$ .009 & .592 $\pm$ .038 \\
BM25          & .533 $\pm$ .037 & .365 $\pm$ .032 & .163 $\pm$ .010 & .556 $\pm$ .043 \\
LSI           & .511 $\pm$ .041 & .351 $\pm$ .036 & .177 $\pm$ .010 & .603 $\pm$ .039 \\
\midrule
Full Agent    & .639 $\pm$ .032 & .436 $\pm$ .030 & .199 $\pm$ .010 & .673 $\pm$ .034 \\
\bottomrule
\end{tabular}}
\end{table}

Even under the legacy protocol, and with ranking now measured independently of the accept/abstain decision, the agent achieves the highest $\text{R}_\mu$@5 (0.673) among the deterministic methods, a 7-point absolute improvement over LSI (0.603), the strongest lexical baseline on this metric.
This is notable because LSI was specifically designed to capture latent semantic relationships~\citep{MarcusMaletic2003}; the agent's reformulation mechanism provides a more targeted form of semantic bridging on this vocabulary-mismatched corpus.


\noindent\textbf{The CSF 1.1 corpus.}
\label{sec:results-real}

Table~\ref{tab:real} reports the same protocol applied to the NIST CSF to SP~800-53r5 corpus.
Absolute numbers are markedly lower across the board, reflecting a genuinely harder task: 300 candidate controls instead of 110, an average of 4.7 gold controls per requirement, and no controls authored to be easy.

\begin{table}[htbp]
\caption{Results on the CSF~1.1 corpus. Repeated-holdout protocol, mean $\pm$ 95\% CI over 30 seeds, 106 requirements, 300 controls, 495 links. The two blocks use different retrieval protocols and are not directly comparable; see the text. Descriptive, as in Table~\ref{tab:ablation}.}
\label{tab:real}
\centering
\setlength{\tabcolsep}{3pt}
\small
\resizebox{\ifdim\width>\columnwidth\columnwidth\else\width\fi}{!}{%
\begin{tabular}{@{}lcccccc@{}}
\toprule
Method & Top-1 & MRR@5 & nDCG@5 & R@5 & Cov. & SelAcc \\
\midrule
\multicolumn{7}{@{}l}{\emph{End-to-end: ranks all 300 controls}} \\
TF-IDF        & .390 $\pm$ .036 & .505 $\pm$ .028 & .307 $\pm$ .022 & .262 $\pm$ .021 & .792 & .423 \\
BM25          & .368 $\pm$ .036 & .489 $\pm$ .027 & .303 $\pm$ .022 & .275 $\pm$ .022 & .778 & .448 \\
LSI           & .360 $\pm$ .033 & .471 $\pm$ .030 & .306 $\pm$ .023 & .269 $\pm$ .022 & .797 & .419 \\
Dual-enc.     & \textbf{.402} $\pm$ .035 & \textbf{.530} $\pm$ .029 & \textbf{.352} $\pm$ .024 & \textbf{.313} $\pm$ .021 & .733 & .489 \\
\midrule
Single        & .368 $\pm$ .036 & .489 $\pm$ .027 & .303 $\pm$ .022 & .275 $\pm$ .022 & .778 & .448 \\
+Multi        & .403 $\pm$ .027 & .516 $\pm$ .023 & .328 $\pm$ .021 & .288 $\pm$ .021 & .746 & .460 \\
+Reform       & .406 $\pm$ .028 & .508 $\pm$ .025 & .317 $\pm$ .022 & .275 $\pm$ .021 & .733 & .489 \\
+Decomp       & .406 $\pm$ .028 & .508 $\pm$ .025 & .317 $\pm$ .022 & .275 $\pm$ .021 & .733 & .489 \\
+CrossRef     & .406 $\pm$ .028 & .508 $\pm$ .025 & .317 $\pm$ .022 & .275 $\pm$ .021 & .733 & .489 \\
Full Agent    & .406 $\pm$ .028 & .508 $\pm$ .025 & .317 $\pm$ .022 & .275 $\pm$ .021 & .721 & .478 \\
\midrule
\multicolumn{7}{@{}l}{\emph{Conditional: reranks the dual-encoder top-20, Recall@20 ceiling .906}} \\
Cross-enc.    & .562 $\pm$ .032 & .666 $\pm$ .026 & .454 $\pm$ .021 & .398 $\pm$ .022 & .775 & .639 \\
\bottomrule
\end{tabular}}
\end{table}

Three results stand out, and they qualify rather than confirm the diagnostic-benchmark findings of Section~\ref{sec:results-synth}.
First, on this corpus no agent component significantly improves ranking.
Neither fusion (multi vs.\ single Top-1 $+0.026$, $p = 0.51$) nor reformulation (Table~\ref{tab:multi}: $+0.028$, $p = 0.38$) produces a significant change.
The seed-averaged means drift upward slightly (single 0.368 to multi 0.403), which a naive across-seed $t$-test would flag as significant, but that apparent effect does not survive the pseudoreplication-safe requirement-level test.
One explanation offers itself immediately: CSF subcategories and SP~800-53 controls are both written in NIST's institutional vocabulary, so there may be less abstraction distance for a thesaurus to bridge.
We record that as a conjecture, not a finding. Our gap measure is fitted per corpus and gives no scale on which one corpus can be called wider-gapped than another, so nothing here tests it, and the next two subsections give reasons to doubt it.
Decomposition and corroboration again leave the ranking untouched (decomposition rarely triggers on single-clause subcategories; a single-framework corpus offers no related frameworks to corroborate against), and verification, a decision-only tool, lowers coverage from 0.733 to 0.721 while nudging selective accuracy, correctly abstaining on some mappings at a coverage cost.
Second, no deterministic configuration out-ranks neural retrieval under the same protocol.
The full agent (0.406 Top-1) is statistically indistinguishable from the strongest lexical baseline (TF-IDF, 0.390; requirement-level $\Delta=+0.016$, $p=0.65$) and from the dual-encoder (0.402; $p=0.78$).
Every method in that comparison ranks all 300 controls.
The cross-encoder does not, which is why Table~\ref{tab:real} places it in a separate block.
It reranks the 20 candidates the dual-encoder returns, so its 0.562 Top-1 is the score of a two-stage pipeline and is capped by that shortlist: the dual-encoder puts a gold control in the top 20 for 90.6\% of requirements, and for the remaining 10 the cross-encoder cannot recover a correct answer at any rank.
That difference is real and we report it as such: as a complete system, dual-encoder retrieval followed by cross-encoder reranking beats every single-stage method here, by $+0.154$ against the agent (95\% CI $[+0.058, +0.250]$, $p=0.002$).
Two readings of that number are worth keeping apart. As an end-to-end comparison between deployable systems it stands, and it is the strongest configuration we measured on this corpus. As a comparison between reranking methods it is bounded above by the shortlist: the cross-encoder can never recover the 9.4\% of requirements whose gold control the first stage missed, and a different first stage would move its score without any change to the reranker. We report the ceiling alongside the score for that reason, not to discount it.
Third, taken with Section~\ref{sec:results-synth}, the deterministic agent never significantly out-ranks a neural baseline under a matched protocol on either benchmark.
It holds no decision-quality edge either: TF-IDF has higher coverage (0.792 vs.\ 0.721) and the cross-encoder pipeline higher selective accuracy (0.639 vs.\ 0.478).

Taken alone, these two corpora invite a tidy conclusion: reformulation helps where the vocabulary gap is wide and is inert where it is narrow. That is what an earlier version of this work concluded, on two corpora. The next three subsections put all four corpora on one protocol and test the reading.

\subsection{Does the reformulation benefit generalize? (RQ1)}
\label{sec:results-multi}

Table~\ref{tab:multi} gives the reformulation contrast on all four corpora under the primary protocol: one scoring pass per requirement over the full control corpus, with LSI fitted on control documents alone. The four figures are produced the same way and can be compared with each other.

\begin{table}[htbp]
\caption{Effect of domain-aware reformulation on Top-1 accuracy, measured identically on all four corpora. One scoring pass per requirement over the full control corpus; LSI fitted on controls only. Intervals are requirement-level bootstrap, $p$ from paired sign-flip permutation. W/L/T counts requirements by whether reformulation turns a wrong top-1 into a right one, a right one into a wrong one, or leaves correctness unchanged. Ties include requirements whose top-ranked control changes without changing correctness.}
\label{tab:multi}
\centering
\setlength{\tabcolsep}{3pt}
\small
\resizebox{\ifdim\width>\columnwidth\columnwidth\else\width\fi}{!}{%
\begin{tabular}{@{}lrrrrrl@{}}
\toprule
Corpus & $n$ & +Multi & +Reform & $\Delta$Top-1 & 95\% CI & W/L/T \\
\midrule
Diagnostic (ours) & 58 & .466 & .586 & $+.121$ & $[+.034, +.224]$ & 8/1/49 \\
\midrule
CSF~1.1 & 106 & .377 & .406 & $+.028$ & $[-.009, +.075]$ & 4/1/101 \\
HIPAA & 68 & .309 & .309 & $\pm.000$ & $[-.059, +.059]$ & 2/2/64 \\
Privacy Fw. & 94 & .340 & .372 & $+.032$ & $[-.011, +.085]$ & 4/1/89 \\
\bottomrule
\end{tabular}}
\end{table}

The estimate on the corpus we built is the only one distinguishable from zero ($p = 0.039$, against $0.376$, $1.000$ and $0.376$), and it is nominal: this analysis is exploratory, the four tests are not corrected as a family, and $p = 0.039$ would not survive a correction across them. We give the four estimates rather than a ratio between them. A ratio would invite the reading that engineered benchmarks inflate effects by some factor, which one corpus of our own against three correlated NIST corpora cannot support (Section~\ref{sec:threats}).

The win/loss/tie counts are the more important column. Top-1 is binary for each requirement, so these means are decided entirely by the requirements on which the two methods disagree, and there are very few of them: nine on the diagnostic corpus, five on CSF, four on HIPAA, five on the Privacy Framework. The $+0.121$ is eight requirements gained and one lost out of 58. We report the counts because an interval on its own does not tell a reader that a corpus-level conclusion rests on this many individual decisions, and because they set the scale on which the sensitivity analysis below should be read.

Reformulation is not idle on the external corpora; it is active and unhelpful, which is a different thing. Four counts separate what it does from what it achieves. The ambiguity gate fires on 58 of 58 diagnostic requirements and on 105/106, 68/68 and 93/94 of the external ones. An expansion is actually produced on 52, 62, 48 and 78, since a fired gate yields nothing when no thesaurus pattern matches. That expansion changes the ranking somewhere on 21, 21, 16 and 35, and moves the top-ranked control on 15, 14, 10 and 18. Only 9, 5, 4 and 5 changes cross the correctness boundary. Most of the reordering on the external corpora moves one wrong control above another, and it is the last of these five numbers, not the first, that drives Top-1.

Comparing the four $p$-values is not itself evidence that the effects differ, and the diagnostic corpus has the smallest $n$, so part of the difference in significance is a difference in power. We therefore quantify the contrast directly: $+0.121$ on the diagnostic corpus ($n = 58$) against $+0.022$ pooled over the external corpora ($n = 268$), a difference of $+0.098$ with a bootstrap 95\% interval of $[+0.001, +0.206]$ \emph{conditional on these four corpora}. The lower end of that interval sits almost exactly on zero, so even taken on its own terms it locates a difference without measuring one.

We stop short of calling that a test of a corpus-type effect. Corpus type is a property of corpora, not of requirements, and this study has four corpora with exactly one of them ours. The effective sample size for any question about corpus type is four clusters, not 326 requirements, and no procedure identifies a corpus-type effect from a single observation in one group. A requirement-level permutation over the same data returns a small $p$-value, but it assumes requirements are exchangeable across corpus types when requirements within a corpus share authors, construction procedure and target catalogue. We report it as a labelled sensitivity analysis and draw no inference from it. Answering the population question needs more corpora, not more requirements inside the ones we have.

The defensible answer to RQ1 is: the effect observed on the corpus we built did not generalize to the three external corpora.

\noindent\textbf{Equivalence, as a secondary result.}
It is tempting to strengthen ``not significant'' into ``equivalent to nothing''. These data support that only weakly. Against the pre-specified $\pm 0.05$ margin, the two one-sided tests give $p_{\text{TOST}} = 0.048$ on HIPAA, $0.152$ on CSF and $0.224$ on the Privacy Framework. Only HIPAA meets the margin, and it does so at $0.048$ against $\alpha = 0.05$, a decision that its four discordant requirements could reverse. We flag it as boundary-sensitive and headline none of it. An earlier version of this work reported CSF and the Privacy Framework as equivalent to no effect; those figures came from the repeated-holdout regime and they do not survive the move to a common protocol. We did not implement an exact matched-pair equivalence procedure; the equivalence results reported here use the stated $t$-based approximation, which is why we treat them as secondary and flag decisions near $\alpha$. The exact sign tests we do report address superiority only.

\noindent\textbf{What this does not establish.}
It does not establish why the split falls where it does. Our corpus differs from the external three in authorship, and also in size, construction procedure, subject breadth and provenance. With one corpus on that side of the comparison these cannot be separated. Section~\ref{sec:threats} returns to it.

\subsection{How much of this depends on choices unrelated to reformulation?}
\label{sec:results-lsi}

Two things changed between the figures an earlier version of this work reported and the figures in Table~\ref{tab:multi}: the evaluation protocol, and the population the latent-semantic backend is fitted on. Neither is about reformulation. We separate them with a three-cell factorial, reported as an unplanned analysis written after the per-corpus effects were known.

\begin{table}[htbp]
\caption{Separating the evaluation protocol from the LSI fitting population. Cell~A is repeated holdouts with LSI fitted on controls only; cell~B is repeated holdouts with LSI fitted on controls plus the training and calibration requirements, which is what the earlier version of this work used; cell~C is the primary one-pass protocol with LSI fitted on controls only. A~vs.~B isolates the fitting population, A~vs.~C isolates the protocol. Unplanned analysis.}
\label{tab:lsi}
\centering
\setlength{\tabcolsep}{3pt}
\small
\resizebox{\ifdim\width>\columnwidth\columnwidth\else\width\fi}{!}{%
\begin{tabular}{@{}lrrrrr@{}}
\toprule
& A & B & C & \multicolumn{2}{c}{attributable to} \\
\cmidrule(l){5-6}
Corpus & hold/ctrl & hold/train & 1-pass/ctrl & LSI fit & protocol \\
\midrule
Diagnostic  & $+.121$ & $+.140$ & $+.121$ & $+.020$ & $.000$ \\
CSF~1.1     & $+.028$ & $+.009$ & $+.028$ & $-.020$ & $.000$ \\
HIPAA       & $\pm.000$ & $+.008$ & $\pm.000$ & $+.008$ & $.000$ \\
Privacy Fw. & $+.032$ & $+.003$ & $+.032$ & $-.029$ & $.000$ \\
\bottomrule
\end{tabular}}
\end{table}

The protocol column is zero on every corpus, to four decimal places. That is close to an identity rather than a discovery, and it is worth stating precisely what it does and does not cover. Every component of this pipeline is deterministic given its inputs, and under a controls-only fit none of those inputs depends on which requirements land in the training split. A requirement therefore scores identically every time it is tested, and the mean over splits equals the single-pass value. The claim is confined to ranking estimates from this pipeline under this fit. It says nothing about evaluation protocols in general, and it does not hold for the decision metrics of Section~\ref{sec:results-openworld}, whose thresholds are calibrated per split. We report it because an earlier draft attributed part of the difference to the protocol, and it does not belong there.

The fitting column does the work, and it has no consistent sign: $+0.020$, $-0.020$, $+0.008$, $-0.029$. Fitting LSI on the training and calibration requirements as well as the controls neither reliably inflates nor reliably deflates the estimate; it moves it by two or three points in whichever direction. A third fitting population, controls plus every requirement including the scored ones, gives $+0.121$, $0.000$, $-0.015$ and $+0.011$, so across the three regimes the CSF estimate spans $0.028$, the HIPAA estimate $0.023$ and the Privacy Framework estimate $0.029$.

That span is as large as the effects themselves on the external corpora. Two things follow. The external estimates are not precise enough to support a claim about the sign of the reformulation effect, quite apart from its significance. And the diagnostic estimate stays between $+0.121$ and $+0.140$ across all three regimes, which we note without explaining: with nine discordant requirements, a point estimate holding still across three fits is weak evidence about anything, and we do not offer it as a property of the corpus.

Cell~B reproduces the four figures the earlier version of this work reported, to four decimal places and including their $p$-values. That is our check that the reimplementation behind this factorial is faithful.

\subsection{Does vocabulary gap identify beneficiaries within a corpus? (RQ2)}
\label{sec:results-gap}

A natural explanation for RQ1's result is that reformulation helps requirements whose vocabulary is distant from their controls, and that our corpus simply contains more of them. The first half of that claim is testable here; the second is not, and we are explicit about which is which.

Our gap measure is computed with a frozen tokenizer from IDF fitted \emph{within each corpus}. That makes gap values comparable between requirements of the same corpus and \textbf{not} comparable between corpora: a gap of $0.8$ in HIPAA and a gap of $0.8$ in CSF are not the same quantity, because the IDF weighting behind them was fitted on different vocabularies. We accordingly standardize gap within corpus before pooling and include corpus fixed effects, which removes between-corpus differences by construction.

The estimand is therefore \emph{within-corpus}: among requirements of the same corpus, do those with relatively larger gaps gain more from reformulation? Fitted to the same one-pass differences as Table~\ref{tab:multi} and pooled over the three external corpora, the interaction between standardized gap and reformulation benefit is $+0.013$, 95\% CI $[-0.011, +0.038]$, sign-flip permutation $p = 0.311$, $n = 268$.

The answer to RQ2 is that we found no evidence of moderation: the interval spans zero and is consistent with a small effect in either direction, so this is a failure to detect one rather than a demonstration that none exists.

Adding our own corpus to the same model moves the coefficient to $+0.021$, 95\% CI $[-0.000, +0.044]$, $p = 0.069$, $n = 326$. We report it because it is informative, and we do not test the difference between the two coefficients. The larger figure is still inconclusive: its interval reaches zero. All we can say is that the point estimate rises when our corpus is included, which is consistent with the moderation being carried by that corpus and equally consistent with noise.

Two further observations bound what the gap measure does. It predicts \emph{lexical failure} reliably: Spearman correlations between gap and TF-IDF correctness are $-.754$, $-.359$, $-.272$ and $-.475$ across the four corpora, all $p < .05$. It does not consistently predict \emph{semantic advantage}: the same correlation against the dual-encoder's margin over TF-IDF gives $+.335$, $-.082$, $-.143$ and $+.302$, significant on two corpora, indistinguishable from zero on two, and inconsistent in sign. Both outcomes here are binary or ternary and therefore heavily tied against a continuous gap score, so these $p$-values come from permuting the gap vector rather than from the asymptotic approximation, and they are Holm-Bonferroni adjusted within each family of four. The pattern is unchanged by the adjustment: all four lexical-failure correlations survive it, and the same two semantic-advantage correlations do. A wide gap tells you term matching will struggle. It does not reliably tell you that a thesaurus will help.

\subsection{What the three answers do and do not license}
\label{sec:results-license}

Taken together, RQ1, the sensitivity analysis and RQ2 support a narrower claim than it is tempting to make.

They support this. Under one common protocol, the reformulation benefit we demonstrated on our own corpus does not appear on three externally authored corpora; each per-corpus estimate turns on fewer than ten requirements; the external estimates move by as much as their own magnitude when a fitting choice unrelated to reformulation is varied; and within the external corpora we found no evidence that a requirement-level lexical gap identifies who benefits.

They do not support a claim that vocabulary-gap severity explains the difference between corpora. Our design cannot address that question at all. Because IDF is corpus-relative we have no common scale on which to say one corpus has a wider gap than another, and the corpus fixed effects remove exactly the between-corpus variation such a claim would need. Any statement of the form ``corpus A has a wider gap than corpus B'' is outside what this measure supports, and we make none.

Nor do they identify authorship as the operative property. The split aligns with our corpus against the NIST-derived ones, but the design cannot separate benchmark authorship from other properties shared across the NIST ecosystem, and one corpus on our side cannot separate authorship from size, subject breadth or construction procedure. The three external corpora are also correlated with one another, sharing a lineage and a target catalogue, so three nulls are weaker evidence than three independent nulls would be.

Nor, finally, is this a result about the thesaurus being poorly built. The mechanism does what it is meant to do on the corpus that rewards it. The point is narrower and less comfortable: we wrote the thesaurus and we wrote the corpus that demonstrates its value, and that corpus did not predict how the thesaurus behaves elsewhere.

This leaves the agent without a demonstrated case on accuracy, and we do not construct one from its other properties. Recording each step establishes where a mapping came from; it does not establish that the record is useful to an auditor, which would need a study with practitioners, and we did not measure execution determinism under pinned dependencies.

\subsection{Open-world gap detection}
\label{sec:results-openworld}

The evaluations above assume every requirement has at least one correct control.
In deployment, the more consequential question is often the open-world one: can the system recognize when \emph{no} adequate control exists and abstain, flagging a compliance gap?
To test this directly, we hold out a random 50\% of controls per seed on the diagnostic benchmark; a test requirement whose gold controls are all removed becomes a constructed no-match query on which the correct behaviour is to abstain. These are manufactured rather than naturally occurring gaps, which Section~\ref{sec:threats} treats as a limit on what the figures below can be read to mean.
This yields on average 4.4 no-match queries per seed (132 across the 30 seeds), a small denominator that we report explicitly and that widens the confidence intervals below.
We report \emph{gap detection}, the fraction of no-match queries on which the system abstains, for three configurations that isolate the decision tools, each calibrated to the same 0.80 target coverage on answerable queries (Table~\ref{tab:openworld}).

\begin{table}[htbp]
\caption{Open-world gap detection on the diagnostic benchmark (50\% of controls held out; mean $\pm$ 95\% CI over 30 seeds; 4.4 no-match queries per seed, 132 total). Gap detection is the abstention rate on constructed no-match queries, made by removing a requirement's gold controls.}
\label{tab:openworld}
\centering
\setlength{\tabcolsep}{3pt}
\small
\resizebox{\ifdim\width>\columnwidth\columnwidth\else\width\fi}{!}{%
\begin{tabular}{@{}lcc@{}}
\toprule
Configuration & Gap detection & Coverage \\
\midrule
Single (conservative)      & .403 $\pm$ .114 & .678 \\
+CrossRef                  & .211 $\pm$ .095 & .828 \\
Full Agent (+Verify)       & .235 $\pm$ .088 & .808 \\
\bottomrule
\end{tabular}}
\end{table}

The results expose a real weakness and, at best, a partial remedy.
The conservative single-backend policy detects 40\% of gaps, but only by keeping overall coverage low (0.68).
Once corroboration and verification are added and the operating point is calibrated toward the 0.80 coverage target, gap detection falls to 0.21--0.24: a system tuned to accept more answerable mappings inevitably accepts more unanswerable ones.
Bidirectional verification provides a small, statistically fragile improvement over corroboration alone (0.211 to 0.235, with overlapping intervals).
Detecting roughly one no-match query in four is far from deployment-ready, and a sensitivity sweep of the verification-tightening constant (Section~\ref{sec:sensitivity}) does not change this regime.
We conclude that abstention calibrated for coverage on answerable queries does not, on its own, solve open-world gap detection, and that a dedicated no-match signal is needed. We state that limitation plainly, and treat it as a priority for future work.

\subsection{Execution traces and sensitivity to the hand-set constants}
\label{sec:sensitivity}

\noindent\textbf{An example trace.}

Each agent decision is accompanied by an execution trace, exported verbatim by the evaluation harness and recording, at the decision step, the effective (post-adjustment) thresholds actually applied.
Table~\ref{tab:trace} reproduces one such trace directly from the exported log (diagnostic benchmark, seed~42), for a GDPR requirement exhibiting vocabulary mismatch.

\begin{table}[htbp]
\caption{Example agent trace, generated directly from the exported log: GDPR Art.~25 (data protection by design and by default), correctly mapped to a privacy-by-design control. Confidence values in Steps~2--4 are fused RRF scores; the reformulation trigger and the effective threshold are logged as shown.}
\label{tab:trace}
\centering
\setlength{\tabcolsep}{3pt}
\small
\resizebox{\ifdim\width>\columnwidth\columnwidth\else\width\fi}{!}{%
\begin{tabular}{@{}clp{0.50\columnwidth}@{}}
\toprule
Step & Action & Observation \\
\midrule
1 & Retrieve & Top: ``Privacy by design checklist integrated into SDLC'' (BM25 conf=10.32) \\
2 & Fuse & Top: ``Data Protection Officer appointed'' (conf=0.049, rel.\ margin=0.01 $<$ 0.10: ambiguous) \\
3 & Reformulate & Expanded with ``privacy by design, PbD, data minimization'' \\
4 & Re-retrieve & Top: ``Privacy by design checklist integrated into SDLC'' (conf=0.049) \\
5 & CrossRef & Corroboration $\rho$=1.0 (domain spans ISO, SOC~2) $\Rightarrow$ relax conf.\ threshold $\times$0.85 \\
6 & Verify & Reverse rank 1/58 within top decile: consistent \\
7 & Decide & \textbf{Accept} (control 8, correct); effective conf.\ threshold 0.045 \\
\bottomrule
\end{tabular}}
\end{table}

The trace shows the mechanism concretely.
Single-shot BM25 (Step~1) happens to rank the correct control first, but fusion (Step~2) reorders a different control to the top with a near-zero relative top-2 margin, which trips the ambiguity test and triggers reformulation.
Reformulation (Step~3) injects the implementation vocabulary, ``privacy by design'' and ``data minimization'', and re-retrieval (Step~4) restores the correct control.
The decision tools then operate without reordering: corroboration finds cross-framework support ($\rho=1.0$) and relaxes the confidence threshold by 15\%, verification confirms bidirectional consistency (the control's text recovers the requirement at reverse rank~1 of 58), and the agent accepts against the logged effective threshold of 0.045.
The example shows how an expansion can move a correct control to the top. On the external corpora the same steps run and the expansion still reorders the ranking, but the reordering rarely moves a correct control into first place (Section~\ref{sec:results-multi}), so the trace looks the same and the outcome does not improve.

\noindent\textbf{The constants.}
The workflow contains four hand-set constants: the relative top-2 margin $\tau_{\text{rel}}$ that triggers reformulation, the two corroboration relaxation factors, and the verification tightening factor.
We check that our conclusions do not depend on their exact values.
We swept $\tau_{\text{rel}}$ over $\{0.05, 0.10, 0.15, 0.20\}$ under the primary protocol, so the output is directly comparable to Table~\ref{tab:multi}. The reformulation effect is identical at every setting on every corpus: $+0.121$, $+0.028$, $0.000$ and $+0.032$, to four decimal places, with a range across settings of exactly zero.
The gate itself is not inert to the constant, which is what makes the invariance informative. Raising $\tau_{\text{rel}}$ from $0.05$ to $0.20$ increases the number of requirements the gate fires on from 100 to 106 of 106 on CSF, from 60 to 68 of 68 on HIPAA, and from 86 to 94 of 94 on the Privacy Framework, and increases the number of expansions correspondingly. Those additional expansions change no outcome. The constant therefore governs how often the tool runs, not what the paper concludes from it, and reformulation's contribution comes from the expansion itself rather than from selective triggering.
An earlier version of this work read the firing rate alone as showing reformulation was always on; Section~\ref{sec:results-multi} gives the four counts that separate invocation from effect.
The decision constants affect only the accept/abstain boundary, and because thresholds are recalibrated to the coverage target, their main measurable effect is on open-world gap detection; sweeping the verification-tightening factor over $\{0.05, 0.10, 0.20\}$ moves gap detection only within $0.25$--$0.31$ at essentially fixed coverage ($\approx 0.79$), so no setting lifts gap detection out of the weak regime reported in Section~\ref{sec:results-openworld}.


%% ================================================================
\section{A preregistered test of a gated hybrid}
\label{sec:hybrid}
%% ================================================================

Everything reported so far is exploratory. The corpora existed before any analysis plan was written, and the plan was not timestamped, so nothing above can claim to have been specified in advance. This section is different, and we describe how it differs before giving the result.

Section~\ref{sec:results} showed that lexical and semantic retrieval fail on different requirements. On the Privacy Framework corpus lexical fusion solves $32/94$ and the dual-encoder $35/94$, yet each uniquely solves requirements the other misses, 14 and 17 respectively, so their oracle union reaches $49/94$. Aggregate ties of this kind conceal complementary errors, and they suggest a combination that is selective rather than uniform: use the dual-encoder when it is confident, and consult the lexical side only when it is not.

That suggestion came from looking at these corpora, which is exactly why the test of it had to be constrained in advance. We therefore froze the design in an executable specification, hashed it, and registered it publicly before computing the contrast.\footnote{\texttt{doi:10.17605/OSF.IO/NZXRV}. The registration records the analysis code, its tests, the input generator, and SHA-256 hashes of every input score matrix.} The registration fixes the primary arm, the comparator, the endpoint, the estimand, the non-inferiority margin, the order of the two tests, the exclusion of the engineered corpus from the primary estimate, and a list of prohibited actions.

We state the limits of that registration plainly. Because the hypothesis was generated from the same corpora that test it, the result is preregistered but still exploratory, and a genuinely confirmatory replication would need a corpus that did not exist when the specification was written. One declared secondary arm, the equal-weight fusion, was computed during method development before registration and is therefore reported as an observed result rather than a preregistered one.

\noindent\textbf{Design.}
The gated arm ranks the full control corpus with the dual-encoder. Only when the dual-encoder's own top-2 relative margin falls below $0.10$ is the lexical side consulted, and the query is then re-ranked by reciprocal rank fusion over TF-IDF, BM25, inductive LSI and the dual-encoder. The gate reads only the dual-encoder scores, so routing is decided before any label is touched. Neither constant is new: $0.10$ is the margin the agent already uses to trigger reformulation (Section~\ref{sec:architecture}), and the five-point margin below was fixed for an earlier, unrelated contrast.

The primary estimand is the mean paired Top-1 difference over the 268 concatenated requirements of the three external corpora, requirement-weighted, with intervals from a bootstrap resampled within corpus. Non-inferiority is tested first at a margin of $0.05$; superiority is tested only if non-inferiority is established, which fixes the multiplicity without a correction.

\noindent\textbf{Result.}
Non-inferiority is established ($p = 0.003$) and superiority is not ($p = 0.48$). The mean difference is $+0.022$, 90\% interval $[-0.022, +0.067]$.

Stated precisely: the analysis rules out mean losses of five Top-1 points or more under the registered model, and provides no evidence of superiority. It does not establish that gating is free. A cost of two or three points is entirely consistent with these data, as the interval shows, and the precision statement registered in advance says as much: at $n = 268$ this design can exclude a five-point loss but cannot resolve anything finer.

That pooled result is the registered answer, and it conceals the more interesting pattern. Table~\ref{tab:hybrid} gives the per-corpus figures, which are descriptive by prior declaration and carry no test.

\begin{table}[htbp]
\caption{Gated hybrid against the dual-encoder alone, Top-1. The pooled contrast is the registered primary result; per-corpus rows are descriptive by prior declaration.}
\label{tab:hybrid}
\centering
\setlength{\tabcolsep}{3pt}
\small
\resizebox{\ifdim\width>\columnwidth\columnwidth\else\width\fi}{!}{%
\begin{tabular}{@{}lccccc@{}}
\toprule
Corpus & $n$ & Dual-enc. & Gated & $\Delta$ & Gate fires \\
\midrule
CSF~1.1 & 106 & .406 & .491 & $+.085$ & .66 \\
HIPAA & 68 & .456 & .368 & $-.088$ & .85 \\
Privacy Fw. & 94 & .372 & .404 & $+.032$ & .74 \\
\midrule
Pooled & 268 & --- & --- & $+.022$ & --- \\
\midrule
\multicolumn{6}{@{}l}{\emph{Engineered corpus, excluded from the pooled estimate}} \\
Diagnostic & 58 & .621 & .569 & $-.052$ & .48 \\
\bottomrule
\end{tabular}}
\end{table}

A gain of $+0.085$ on CSF sits against a loss of $-0.088$ on HIPAA. A study evaluating this gate on either corpus alone would have reported a clear success or a clear harm. The pooled figure averages those two against a third, $+0.032$ on the Privacy Framework, over all 268 requirements. That is the same cross-corpus spread as the reformulation result of Section~\ref{sec:results-multi}, reached here inside a single frozen design rather than assembled after the fact.

The complementarity counts show the mechanism. On CSF the gate recovers 7 of the 11 requirements that lexical fusion uniquely solves while losing 3 of the 14 the dual-encoder uniquely solves. On HIPAA the trade reverses: 2 recovered against 8 lost. The gate is not behaving differently on the two corpora; the corpora differ in whether the lexical side has anything to contribute when the dual-encoder hesitates.

Two features of the data qualify the design itself. The gate fires on 66\% to 85\% of queries, so a $0.10$ relative margin is not selective on this task and the gated arm sits close to the equal-weight fusion, most visibly on HIPAA where it fires on 85\% of queries and the two arms coincide. And the pre-declared hard-fallback variant, which hands gated queries to lexical fusion outright rather than fusing all four backends, is worse than the primary gate on every corpus. Declaring it secondary in advance is what prevents that from being presented as a design choice made on principle.

A sensitivity arm replacing inductive LSI with a transductive fit gives $+0.026$ against the primary's $+0.022$, so the conclusion does not depend on that choice. This arm fits on the controls plus every requirement, including the ones being scored; it is not the fitting population used by the repeated-holdout runs, which excludes the test requirements. The three populations are set out in Section~\ref{sec:results-lsi}.

%% ================================================================
\section{Discussion}
\label{sec:discussion}
%% ================================================================

\subsection{What the four corpora buy that two could not}

With two corpora we had a clean story: reformulation helps where the vocabulary gap is wide and is inert where it is narrow. The mechanism was plausible. Regulatory text is written at a policy abstraction level (``implement appropriate technical measures'') while controls describe specific implementations (``AES-256 with FIPS 140-2''), and a curated thesaurus should bridge that because the mapping between regulatory concepts and implementation terms is relatively stable within a domain.

Two corpora cannot distinguish that explanation from several others. Our corpus differed from the NIST corpus in apparent gap severity, but it also differed in who wrote it, how large it was, and how it was constructed, and with a single corpus on each side these are perfectly confounded. Adding HIPAA and the Privacy Framework does not resolve the confound, but it does remove the gap explanation's most testable form. HIPAA is the one corpus here whose two sides come from different institutions, with statutory text on the requirement side, which is qualitatively where a regulation-to-implementation thesaurus should help most. Its effect is $0.000$, and it sits with the Privacy Framework's $+0.032$ and CSF's $+0.028$ rather than with ours.

We are careful about what follows. We cannot say that HIPAA's measured gap is wider than any other corpus here, because our gap measure is fitted per corpus and provides no common scale. We can say that the corpus with the strongest qualitative reason to benefit did not benefit.

A third corpus and a fourth also changed what a single number is worth. With two corpora the reformulation figure looked like a property of the method under two conditions. With four, and with the fitting sensitivity of Section~\ref{sec:results-lsi} measured, the external figures look like what they are: differences of four or five requirements, movable by roughly their own magnitude by a choice about how to fit a backend.

The lesson we take is not about thesauri, and it is not a demonstration that authorship is the operative variable. It is that a conditional recommendation of the form ``invest in domain glossaries when the vocabulary gap is wide'' was available from two corpora, looked well supported, and does not survive two more. Whatever distinguishes our corpus from NIST's, it was invisible while we had one corpus of each kind.

The thesaurus approach does expose the terms driving each expansion, which a neural reranker~\citep{Nogueira2019} does not.
Whether that inspectability changes an auditor's work is untested here, so we record it as a property of the method rather than as a measured advantage.

\subsection{What a registered null is worth}

The gated hybrid in Section~\ref{sec:hybrid} returned a non-inferiority result without evidence of superiority. Reported without a registration, that result would be hard to weigh, because a reader has no way to know whether it was the analysis we intended in advance or the most favourable of several we could have run. Registration does not make such a result publishable by itself; what it does is constrain outcome-dependent analysis, so that the reported test is demonstrably the one specified in advance. The mechanism is worth spelling out because it is cheap and rarely used in this literature.

Three provisions did the work. The endpoint, margin and test order were fixed in code before the registered contrast was computed, so the superiority test could only run if non-inferiority passed, and no second test could be substituted if the first disappointed. The secondary variants, including the hard fallback that turned out to be worse everywhere, were declared in advance, so the better-performing gate cannot be presented as a principled choice made ahead of time. And the analysis inputs were hashed, so a reader can check that the numbers came from the matrices we said they came from.

The benefit is that a result of this shape, which offers no headline improvement, can be reported as evidence rather than discarded, which matters in a subfield where the space of plausible retrieval combinations is large and the incentive to report only the combination that worked is correspondingly strong.

We also disclose what registration did not fix. The hypothesis was generated from the same corpora that tested it, so the result is preregistered but exploratory, and one secondary arm was computed before registration and is reported as observed rather than preregistered.

\subsection{Ranking, decisions, and what the agent does not buy}

A recurring theme in our results is the distinction between \emph{ranking quality} (how well are controls ordered?) and \emph{decision quality} (when should the system commit to a mapping, and when abstain?).
Ranking is computed from the full candidate list regardless of the accept/abstain decision, and coverage, selective accuracy and open-world gap detection are reported as separate decision metrics. Computing ranking only over accepted queries lets a coverage change masquerade as a ranking gain, which is what an earlier version of this work did.
This separation is what makes the component analysis trustworthy, and it changes the conclusions: several apparent ranking gains dissolve once decoupled from coverage, and once we also apply requirement-level rather than seed-level inference (Section~\ref{sec:protocol-stats}), the agent turns out to hold neither a ranking edge nor a decision-quality edge over the strongest baselines. What remains is a negative result about the components, not a qualitative advantage recovered in their place.

The distinction is nonetheless underappreciated in the traceability literature, which focuses almost exclusively on ranking metrics.
We argue that for deployment in compliance settings, decision quality is at least as important as ranking quality, because incorrect mappings in an accepted state carry real regulatory risk, and because, as Section~\ref{sec:results-openworld} shows, recognizing the \emph{absence} of a valid mapping is itself an unsolved problem.
The decision tools illustrate the point in both directions: corroboration and verification barely move the ranking yet substantially reshape coverage and gap detection, and verification, whose rationale is decision hygiene, must be gated carefully because in densely linked corpora it trades coverage for only a fragile calibration gain.

Once ranking is measured independently of the decision and compared with a requirement-level test, the picture is consistent: the deterministic agent does not out-rank neural retrieval on any corpus.
Under the primary protocol full RAA never exceeds the dual-encoder on any corpus: $-0.034$, $0.000$, $-0.147$ and $0.000$ (Table~\ref{tab:unified}). The repeated-holdout tables agree on the diagnostic corpus, where the two are statistically indistinguishable (0.672 against 0.644 Top-1, $p\approx1.0$), and on the external corpora no deterministic component improves ranking significantly.
A two-stage cross-encoder pipeline scores higher on CSF and the Privacy Framework under the primary protocol ($.547$ against $.406$, and $.457$ against $.372$), and lower than the dual-encoder alone on the diagnostic corpus and HIPAA. It reranks a 20-candidate shortlist, so as a comparison between reranking methods it is bounded by that shortlist and we give the ceiling alongside it.

We draw two conclusions.
First, general-purpose learned representations are strong retrievers for compliance text, and a deterministic pipeline built from lexical backends and a curated thesaurus does not beat them on ranking, even where the thesaurus helps most.
Second, the agent holds neither the best ranking nor the best coverage nor the best selective accuracy on any corpus, and we do not offer its other properties as a substitute.
It does record each step and its inputs, and the neural rerankers do not. That establishes where a mapping came from. It does not establish that the record is useful under audit, which is a claim about practitioners and would need to be tested on them.

The natural next step, combining lexical and semantic retrieval, is the one we registered and tested, and Section~\ref{sec:hybrid} reports what happened: the registered analysis rules out a mean loss of five Top-1 points or more and finds no evidence of superiority, while producing a large gain on one corpus and a nearly equal loss on another. That per-corpus divergence is the same pattern as the reformulation result, arrived at by a different route, and it suggests the useful question is not which combination wins but which corpus properties decide the answer. We do not know what those properties are. A requirement-level lexical gap was the obvious candidate, and Section~\ref{sec:results-gap} finds it does not identify beneficiaries within a corpus.


%% ================================================================
\section{Threats to validity}
\label{sec:threats}
%% ================================================================

\noindent\textbf{Construct validity.}
The diagnostic benchmark is synthetic, and its vocabulary-mismatched controls, though authored to share no content words with their requirements, may not represent the full diversity of real regulatory--implementation gaps.
The three external corpora mitigate this threat for the texts and ground truth themselves, but they capture one family of mapping regimes. All three map into a single NIST control catalogue, and on CSF and the Privacy Framework both sides are NIST products. We do not claim their vocabulary gap is narrow; our measure is corpus-relative and cannot support a statement of that form. What we can say is that, as noted in Section~\ref{sec:dataset-real}, the crosswalks encode NIST's concept-level judgment of correspondence~\citep{NISTIR8477} rather than evidence that an organization has implemented a control satisfying a requirement.
Our claims on those corpora are therefore about reference-link recovery, not implemented-compliance verification.
No corpus here reflects ambiguous multi-clause requirements, implicit dependencies between controls, or organizational context.
Finally, the open-world evaluation manufactures no-match queries by holding out controls; genuine compliance gaps may differ systematically from queries whose gold controls happen to be removed, so the gap-detection figures should be read as a controlled stress test rather than a field estimate or a bound in either direction.

\noindent\textbf{Internal validity.}
The domain thesaurus and concept patterns are hand-crafted by the authors, who are practicing cybersecurity specialists, creating a risk of overfitting the thesaurus to the diagnostic benchmark.
We separated thesaurus design (general domain knowledge) from benchmark design (specific regulatory texts), but that is a procedural safeguard and not a test. Repeated splits do not address it either: resampling the same corpus cannot reveal that the corpus and the thesaurus were built by the same authors.
The external corpora provide the stronger safeguard: the thesaurus was frozen before any of them was assembled, and neither their texts nor their links were authored by us.
The smaller effects on those corpora are therefore not an artifact of tuning against them. They are not a bound on transfer either. Each rests on four or five discordant requirements and moves by roughly its own magnitude under a change of LSI fitting population, so they establish that the diagnostic result did not recur rather than measuring how much of it survives.

\noindent\textbf{External validity: the three external corpora are not independent of one another.}
This is the most important limitation of the study and it bears directly on the central claim.
CSF~1.1 and the Privacy Framework are NIST products. The HIPAA Security Rule is not: it is statutory text issued by the U.S.\ Department of Health and Human Services, and what NIST supplies there is the crosswalk into SP~800-53. All three nonetheless map into the same 300-control corpus, and all three crosswalks come from NIST, two of them from the same programme, so the mappings share a lineage even where the requirement text does not.
They are correlated evaluation corpora, not independent samples of regulatory text.
A null that holds across all three is therefore weaker evidence than a null across three unrelated corpora would be, and we cannot exclude that some shared property of the NIST ecosystem, rather than external authorship as such, is what distinguishes them from the corpus we built.
Testing that requires externally authored corpora from different institutions, ideally with a different target catalogue, and it is the single most valuable extension of this work.

For the same reason we make no quantitative claim about how much engineered benchmarks inflate effects.
The comparison here is one corpus of our own against three correlated external ones. It supports a claim that the initial estimate did not generalize across those three, and not a ratio, and not a formal replication result: the initial estimate was never published, and three corpora sharing a lineage are not three independent attempts.

Our conclusions are also bounded by the set of methods we measured. We compare lexical backends, their fusion, a curated thesaurus, a dual-encoder and a cross-encoder reranker, and we do not evaluate a prompted or instruction-tuned language model as a reranker. Recent work in this area reports that such models can be competitive with supervised rerankers without task-specific training~\citep{Sun2023}, and one of the closest studies to ours compares prompting against conventional traceability methods on regulatory text~\citep{Etezadi2026}. A statement of the form ``no method we tried helps on the external corpora'' therefore does not extend to that family, and the comparison is a natural next step rather than one this paper settles.

The diagnostic benchmark covers 7 frameworks and the external corpora three framework pairs; real organizations may need to map across dozens of standards (e.g., FedRAMP, CMMC, DORA, NIS2), and the thesaurus would require extension for new domains.
Corpus sizes (58, 106, 68 and 94 requirements) limit statistical power, and we do not claim that repeated seeds overcome this: the 30 stratified splits overlap heavily and are not independent replicates, so we base all inference on requirement-level paired tests (Section~\ref{sec:protocol-stats}) rather than on across-seed variance, and we report the small open-world denominator (4.4 no-match queries per seed) explicitly.
Power is not the only limit here, and it may not be the binding one. Each per-corpus reformulation estimate is decided by fewer than ten requirements whose top-1 correctness changes at all (Table~\ref{tab:multi}), and each moves by up to $0.029$ when the latent-semantic backend is refitted on a different population (Table~\ref{tab:lsi}). On the external corpora that movement is comparable to the estimate itself. Readers should treat the external point estimates as locating a small effect rather than as measuring its size or fixing its sign.
The registered hybrid test states its own power limit explicitly: at $n = 268$ the one-sided half-width is $0.048$ against a margin of $0.05$, so it can exclude a mean loss of five points or more but cannot resolve anything finer, and a failure to establish non-inferiority would have been consistent with a power limit rather than with harm.

\noindent\textbf{Statistical conclusion validity.}
Four limitations of our own analysis should be read alongside the results.

The equivalence tests are the weakest inference in the paper. Top-1 is binary per requirement, so the paired differences take three values and are far from normal, and the two one-sided $t$ tests behind $p_{\text{TOST}}$ are an approximation whose quality we have not established at these sample sizes. We report equivalence as a secondary result, flag any decision within $0.01$ of $\alpha$, and headline none of it. We did not implement an exact matched-pair equivalence procedure; exact procedures for paired binary outcomes exist, and their absence here is a limitation of this analysis. The exact sign tests reported alongside address superiority only, and an earlier version of this work presented one of them in a way that implied otherwise.

The gap measure is fitted with corpus-relative IDF, so gap values are comparable within a corpus but not across corpora. All pooled analyses standardize within corpus before pooling for that reason, and we make no claim that any corpus has a larger absolute gap than another on a common scale.

Section~\ref{sec:results-lsi} and the engineered-versus-external contrast were written after the per-corpus effects were known. They are unplanned analyses, we label them as such wherever they appear, and only Section~\ref{sec:hybrid} carries a preregistration.

The decision metrics come from repeated stratified holdouts and are computed on per-requirement means over an uneven number of splits, so per-requirement variances are unequal there. We report those metrics descriptively and rest no significance claim on them.


%% ================================================================
\section{Conclusion and future work}
\label{sec:conclusion}
%% ================================================================

We evaluated one frozen set of retrieval methods on four corpora of regulatory traceability: a diagnostic instrument we built, and three whose requirement texts are externally authored framework or statutory documents and whose links are NIST-authored crosswalks.
The study began as an attempt to characterize when a curated compliance thesaurus pays off. It ended as a study of what one benchmark can establish about a method.

Four findings follow.

First, the $+0.121$ estimate on the corpus we built ($p = 0.039$) did not recur under a common protocol. The three external corpora give $+0.028$, $0.000$ and $+0.032$, none distinguishable from zero, and all three remain imprecise for the reasons in the next paragraph. An earlier version of this work, with our corpus and one external corpus, concluded that the benefit was real but conditional on vocabulary-gap width. Two further corpora do not support that reading.

Second, the estimates are less solid than four decimal places suggest. Each turns on fewer than ten requirements whose top-1 correctness changes at all, and each moves by up to $0.029$ when the latent-semantic backend is refitted on a different population, a choice with nothing to do with reformulation. On the external corpora that movement is the size of the effect. The evaluation protocol, which we had previously suspected, contributes nothing at all: repeated holdouts and a single full-population pass agree exactly, because with a controls-only fit nothing in the pipeline depends on the split.

Third, we found no evidence that requirement-level vocabulary distance identifies who benefits within a corpus. Pooled over the external corpora the interaction between standardized gap and reformulation benefit is $+0.013$ (95\% CI $[-0.011, +0.038]$, $p = 0.31$), which is inconclusive: the interval is consistent with a small effect in either direction and we do not read it as a null. Gap predicts lexical failure reliably and semantic advantage inconsistently. We cannot go further and say that gap explains the difference \emph{between} corpora: the measure is fitted per corpus and so provides no common scale, and with one corpus of our own the design cannot separate authorship from size, construction or subject breadth. What we can say is that we wrote both the thesaurus and the corpus on which the thesaurus succeeds, and that this corpus did not predict how the thesaurus behaves elsewhere.

Fourth, a preregistered test of a gated lexical-semantic hybrid rules out a mean loss of five Top-1 points or more against semantic retrieval alone ($p = 0.003$) and finds no evidence of superiority ($p = 0.48$), while producing $+0.085$ on one corpus against $-0.088$ on another. A cost of two or three points remains consistent with those data. The design was frozen, hashed and publicly registered before it ran, which constrains outcome-dependent analysis and lets a result of this shape be reported as evidence rather than set aside.

The negative results about our own system stand: no added deterministic RAA component significantly improves ranking on any external corpus, the agent holds no coverage or selective-accuracy advantage, and its calibrated abstention detects only a minority of constructed no-match cases, so it is not ready for gap discovery. We do not substitute an untested claim about auditability for the missing metric case.

The most valuable extension is the one this study cannot perform on itself. Our three external corpora share a NIST lineage and a common target catalogue, so they are correlated evidence rather than three independent attempts, and distinguishing external authorship from the properties of one institution's document ecosystem requires corpora from different institutions with a different control catalogue. Beyond that: designing a dedicated no-match signal to make open-world gap detection reliable; admitting domain-adapted encoders such as Legal-BERT~\citep{Chalkidis2020} or SecureBERT~\citep{Aghaei2022} as backends; and investigating temporal drift, where regulations and controls evolve and previously valid mappings become stale~\citep{Gama2014}.

A closing methodological point. The reversal reported here was possible only because the same frozen methods were run on more than one externally authored corpus, and because we then checked how much the surviving numbers depended on choices we had not thought of as choices. Had we stopped at two corpora, we would have submitted a plausible conditional recommendation that our own later data contradicts. In a subfield where public collections are scarce and authors routinely build their own, a single benchmark, however carefully constructed, is not sufficient evidence that a method transfers, and a point estimate is worth little without the count of decisions behind it.


%% ================================================================
%% CRediT and Declarations
%% ================================================================


\section*{Statements and Declarations}

\subsection*{Funding}
The authors received no funding for this work.

\subsection*{Competing interests}
The authors declare that they have no known competing financial interests or
personal relationships that could have appeared to influence the work
reported in this paper.

\subsection*{Ethics approval and consent to participate}
Not applicable. This study involves no human participants, human data or
animals.

\subsection*{Consent for publication}
Not applicable.

\subsection*{Data availability}
All four evaluation corpora accompany this article as Online Resource~1: the
diagnostic benchmark (58 requirements, 110 controls, 86 links), and the
CSF~1.1 (106 subcategories, 495 links), HIPAA Security Rule (68 requirements,
274 links) and Privacy Framework (94 subcategories, 456 links) corpora, which
share one 300-control corpus drawn from SP~800-53r5. Online Resource~1 also
carries the provenance record written by each corpus builder, the
reproduction guide and the source manifest. The scripts that regenerate the
three external corpora from published federal source texts, the NIST OSCAL
catalogue and NIST crosswalk artifacts are in the code repository below. The
corpora are licensed CC BY 4.0; the underlying federal source artifacts are
United States Government works in the public domain.

\subsection*{Materials and code availability}
The evaluation framework, all baselines, the agent implementation, the domain
thesaurus and concept patterns, the ablation and open-world harness, the
analysis scripts and the trace-export utility are available at
\url{https://github.com/sumanshusohal/RAA-Compliance-Mapping} under the MIT
Licence. The preregistration of the gated hybrid analysis is at
\url{https://doi.org/10.17605/OSF.IO/NZXRV}.

Experiments use recorded inputs, fixed seeds and fixed inference settings,
with Hugging Face model revisions pinned in every path that loads one. Exact
execution stability across hardware and software environments was not
evaluated, so we claim repeatability of our own runs rather than bitwise
reproducibility on arbitrary machines. The 30-seed logs underlying the
holdout tables are included with the release; the one-pass, factorial and
preregistered hybrid results are backed by JSON records carrying their
arguments, git commit, library versions and SHA-256 hashes of every analysis
module and input. \texttt{audit\_records.py} checks every record against the
working tree and against the commit it names, and \texttt{make\_tables.py
--check} verifies that the four primary quantitative tables in this paper
match those records.

\subsection*{Author contributions}
\textbf{Sumanshu Sohal}: conceptualization, methodology, software,
validation, formal analysis, investigation, data curation, writing --
original draft, visualization. \textbf{Darshankumar Prajapati}: methodology,
investigation, validation, writing -- review and editing.

\subsection*{Use of generative AI}
The authors' use of generative AI and AI-assisted tools is documented in
Section~\ref{sec:protocol}, as the journal requires. The authors
independently verified all analyses, statistics, citations, code and results,
reviewed and edited any AI-assisted output, and take full responsibility for
the content of the publication.


\appendix
\section{Domain thesaurus}
\label{app:thesaurus}

Table~\ref{tab:thesaurus} reproduces the complete 20-family compliance thesaurus used by the reformulation tool.
The 32 pattern-based concept-extraction rules are included in the released implementation.

\begin{table}[htbp]
\caption{The 20 concept families of the domain thesaurus.}
\label{tab:thesaurus}
\centering
\setlength{\tabcolsep}{3pt}
\small
\resizebox{\ifdim\width>\columnwidth\columnwidth\else\width\fi}{!}{%
\begin{tabular}{@{}lp{0.70\columnwidth}@{}}
\toprule
Concept family & Expansion terms \\
\midrule
encryption & encrypt, cryptographic, cipher, AES, TLS, SSL, cryptography, encoded \\
access control & authorization, authentication, RBAC, role-based, permissions, privilege, IAM, identity \\
authentication & MFA, multi-factor, biometric, credential, identity verification, login, SSO, password \\
audit & logging, log, audit trail, SIEM, monitoring, record, tracking, accountability \\
monitoring & detection, surveillance, IDS, IPS, EDR, alerting, observability, anomaly detection \\
breach & incident, security event, compromise, violation, data leak, exposure, notification \\
backup & recovery, restoration, disaster recovery, business continuity, redundancy, replication \\
firewall & network boundary, perimeter, DMZ, segmentation, packet filtering, network protection \\
vulnerability & weakness, exposure, CVE, patch, remediation, scanning, penetration test \\
data protection & privacy, confidentiality, PII, personal data, data subject, GDPR, data handling \\
configuration & baseline, hardening, CIS benchmark, IaC, infrastructure as code, system setup \\
incident response & incident handling, IR, playbook, escalation, containment, eradication, recovery \\
risk assessment & risk analysis, threat assessment, risk management, impact analysis, likelihood \\
training & awareness, education, security training, phishing simulation, security culture \\
change management & change control, CAB, change advisory, approval, release management \\
disposal & sanitization, media destruction, data wiping, degaussing, secure deletion, NIST 800-88 \\
transmission & in-transit, data in motion, transport, TLS, encryption in transit, network security \\
supplier & third-party, vendor, supply chain, outsourcing, contractor, processor \\
erasure & deletion, right to be forgotten, data removal, purge, expunge, destroy \\
impact assessment & DPIA, PIA, privacy impact, risk evaluation, data protection impact \\
\bottomrule
\end{tabular}}
\end{table}

\begin{acknowledgements}
We thank the reviewers of the earlier version of this work, whose criticism
prompted the multi-corpus design reported here.


%% ================================================================
%% Bibliography
%% ================================================================
\end{acknowledgements}
\begin{thebibliography}{99}

\bibitem[Abualhaija et~al.(2022)]{Abualhaija2022}
Abualhaija, S., Arora, C., Briand, L.C., 2022. COREQQA: a COmpliance REQuirements understanding using question answering tool. In: Proc.\ ESEC/FSE (Tool Demonstrations). pp.\ 1682--1686. \url{https://doi.org/10.1145/3540250.3558926}

\bibitem[Aghaei et~al.(2022)]{Aghaei2022}
Aghaei, E., Niu, X., Shadid, W., Al-Shaer, E., 2022. SecureBERT: A domain-specific language model for cybersecurity. In: Proc.\ SecureComm. arXiv:2204.02685. \url{https://doi.org/10.1007/978-3-031-25538-0_3}

\bibitem[Akhigbe et~al.(2019)]{Akhigbe2019}
Akhigbe, O., Amyot, D., Richards, G., 2019. A systematic literature mapping of goal and non-goal modelling methods for legal and regulatory compliance. Requir.\ Eng.\ 24~(4), 459--481. \url{https://doi.org/10.1007/s00766-018-0294-1}

\bibitem[Ali et~al.(2013)]{Ali2013}
Ali, N., Gueheneuc, Y.-G., Antoniol, G., 2013. Trustrace: Mining software repositories to improve the accuracy of requirement traceability links. IEEE Trans.\ Softw.\ Eng.\ 39~(5), 725--741. \url{https://doi.org/10.1109/tse.2012.71}

\bibitem[Amaral et~al.(2022)]{Amaral2022}
Amaral, O., Abualhaija, S., Torre, D., Sabetzadeh, M., Briand, L.C., 2022. AI-enabled automation for completeness checking of privacy policies. IEEE Trans.\ Softw.\ Eng.\ 48~(11), 4647--4674. \url{https://doi.org/10.1109/tse.2021.3124332}

\bibitem[Amaral Cejas et~al.(2023)]{Cejas2023}
Amaral Cejas, O., Azeem, M.I., Abualhaija, S., Briand, L.C., 2023. NLP-based automated compliance checking of data processing agreements against GDPR. IEEE Trans.\ Softw.\ Eng.\ 49~(9), 4282--4303. \url{https://doi.org/10.1109/TSE.2023.3288901}

\bibitem[Antoniol et~al.(2002)]{Antoniol2002}
Antoniol, G., Canfora, G., Casazza, G., De~Lucia, A., Merlo, E., 2002. Recovering traceability links between code and documentation. IEEE Trans.\ Softw.\ Eng.\ 28~(10), 970--983. \url{https://doi.org/10.1109/tse.2002.1041053}

\bibitem[Armstrong et~al.(2009)]{Armstrong2009}
Armstrong, T.G., Moffat, A., Webber, W., Zobel, J., 2009. Improvements that don't add up: ad-hoc retrieval results since 1998. In: Proc.\ 18th ACM Conf.\ on Information and Knowledge Management (CIKM). pp.\ 601--610. \url{https://doi.org/10.1145/1645953.1646031}

\bibitem[Aronson and Lang(2010)]{Aronson2010}
Aronson, A.R., Lang, F.-M., 2010. An overview of MetaMap: historical perspective and recent advances. J.\ Am.\ Med.\ Inform.\ Assoc.\ 17~(3), 229--236. \url{https://doi.org/10.1136/jamia.2009.002733}

\bibitem[Borg et~al.(2014)]{Borg2014}
Borg, M., Runeson, P., Ard\"o, A., 2014. Recovering from a decade: a systematic mapping of information retrieval approaches to software traceability. Empir.\ Softw.\ Eng.\ 19~(6), 1565--1616. \url{https://doi.org/10.1007/s10664-013-9255-y}

\bibitem[Bowman and Dahl(2021)]{BowmanDahl2021}
Bowman, S.R., Dahl, G., 2021. What will it take to fix benchmarking in natural language understanding? In: Proc.\ NAACL-HLT. pp.\ 4843--4855. \url{https://doi.org/10.18653/v1/2021.naacl-main.385}

\bibitem[Bran et~al.(2024)]{Bran2024}
Bran, A.M., Cox, S., Schilter, O., Baldassari, C., White, A.D., Schwaller, P., 2024. Augmenting large language models with chemistry tools. Nat.\ Mach.\ Intell.\ 6, 525--535. \url{https://doi.org/10.1038/s42256-024-00832-8}

\bibitem[Breaux and Ant\'on(2008)]{Breaux2008}
Breaux, T.D., Ant\'on, A.I., 2008. Analyzing regulatory rules for privacy and security requirements. IEEE Trans.\ Softw.\ Eng.\ 34~(1), 5--20. \url{https://doi.org/10.1109/tse.2007.70746}

\bibitem[Carpineto and Romano(2012)]{Carpineto2012}
Carpineto, C., Romano, G., 2012. A survey of automatic query expansion in information retrieval. ACM Comput.\ Surv.\ 44~(1), 1:1--1:50. \url{https://doi.org/10.1145/2071389.2071390}

\bibitem[Chalkidis et~al.(2020)]{Chalkidis2020}
Chalkidis, I., Fergadiotis, M., Malakasiotis, P., Aletras, N., Androutsopoulos, I., 2020. LEGAL-BERT: The Muppets straight out of Law School. In: Findings of EMNLP. pp.\ 2898--2904. \url{https://doi.org/10.18653/v1/2020.findings-emnlp.261}

\bibitem[Cleland-Huang et~al.(2014)]{Cleland-Huang2014}
Cleland-Huang, J., Gotel, O.C.Z., Huffman~Hayes, J., Mader, P., Zisman, A., 2014. Software traceability: trends and future directions. In: Proc.\ Future of Software Engineering (FOSE). pp.\ 55--69. \url{https://doi.org/10.1145/2593882.2593891}

\bibitem[Cleland-Huang et~al.(2010)]{ClelandHuang2010}
Cleland-Huang, J., Czauderna, A., Gibiec, M., Emenecker, J., 2010. A machine learning approach for tracing regulatory codes to product specific requirements. In: Proc.\ 32nd ACM/IEEE ICSE. pp.\ 155--164. \url{https://doi.org/10.1145/1806799.1806825}

\bibitem[Cormack et~al.(2009)]{Cormack2009}
Cormack, G.V., Clarke, C.L.A., Buettcher, S., 2009. Reciprocal rank fusion outperforms condorcet and individual rank learning methods. In: Proc.\ 32nd ACM SIGIR. pp.\ 758--759. \url{https://doi.org/10.1145/1571941.1572114}

\bibitem[Cuddeback et~al.(2010)]{Cuddeback2010}
Cuddeback, D., Dekhtyar, A., Huffman~Hayes, J., 2010. Automated requirements traceability: the study of human analysts. In: Proc.\ 18th IEEE RE. pp.\ 231--240. \url{https://doi.org/10.1109/re.2010.35}

\bibitem[Devlin et~al.(2019)]{Devlin2019}
Devlin, J., Chang, M.-W., Lee, K., Toutanova, K., 2019. BERT: Pre-training of deep bidirectional transformers for language understanding. In: Proc.\ NAACL-HLT. pp.\ 4171--4186. \url{https://doi.org/10.18653/v1/N19-1423}

\bibitem[El-Yaniv and Wiener(2010)]{ElYaniv2010}
El-Yaniv, R., Wiener, Y., 2010. On the foundations of noise-free selective classification. J.\ Mach.\ Learn.\ Res.\ 11, 1605--1641.

\bibitem[Etezadi et~al.(2026)]{Etezadi2026}
Etezadi, R., Abualhaija, S., Arora, C., Briand, L.C., 2026. Classifier or prompt: a case study on legal requirements traceability. Empir.\ Softw.\ Eng.\ 31, 85. \url{https://doi.org/10.1007/s10664-026-10827-1}

\bibitem[Fox and Shaw(1994)]{Fox1994}
Fox, E.A., Shaw, J.A., 1994. Combination of multiple searches. In: Proc.\ 2nd TREC. pp.\ 243--252. \url{https://doi.org/10.6028/nist.sp.500-215.vt}

\bibitem[Gama et~al.(2014)]{Gama2014}
Gama, J., \v{Z}liobait\.{e}, I., Bifet, A., Pechenizkiy, M., Bouchachia, A., 2014. A survey on concept drift adaptation. ACM Comput.\ Surv.\ 46~(4), 44:1--44:37. \url{https://doi.org/10.1145/2523813}

\bibitem[Geifman and El-Yaniv(2017)]{Geifman2017}
Geifman, Y., El-Yaniv, R., 2017. Selective classification for deep neural networks. In: Proc.\ NeurIPS. pp.\ 4878--4887. \url{https://doi.org/10.48550/arXiv.1705.08500}

\bibitem[Geifman and El-Yaniv(2019)]{Geifman2019}
Geifman, Y., El-Yaniv, R., 2019. SelectiveNet: A deep neural network with an integrated reject option. In: Proc.\ ICML. pp.\ 2151--2159. \url{https://doi.org/10.48550/arXiv.1901.09192}

\bibitem[Ghanavati et~al.(2014)]{Ghanavati2014}
Ghanavati, S., Amyot, D., Rifaut, A., 2014. Legal goal-oriented requirement language (Legal GRL) for modeling regulations. In: Proc.\ 6th Int.\ Workshop on Modeling in Software Engineering (MiSE). pp.\ 1--6. \url{https://doi.org/10.1145/2593770.2593780}

\bibitem[Guo et~al.(2017)]{Guo2017}
Guo, J., Cheng, J., Cleland-Huang, J., 2017. Semantically enhanced software traceability using deep learning techniques. In: Proc.\ 39th ICSE. pp.\ 3--14. \url{https://doi.org/10.1109/icse.2017.9}

\bibitem[Hashmi et~al.(2018)]{Hashmi2018}
Hashmi, M., Governatori, G., Lam, H.-P., Wynn, M.T., 2018. Are we done with business process compliance: state of the art and challenges ahead. Knowl.\ Inf.\ Syst.\ 57~(1), 79--133. \url{https://doi.org/10.1007/s10115-017-1142-1}

\bibitem[Ingolfo et~al.(2011)]{Ingolfo2011}
Ingolfo, S., Siena, A., Mylopoulos, J., 2011. Establishing regulatory compliance for software requirements. In: Proc.\ 30th Int.\ Conf.\ on Conceptual Modeling (ER). pp.\ 47--61. \url{https://doi.org/10.1007/978-3-642-24606-7_5}

\bibitem[Kamath et~al.(2020)]{Kamath2020}
Kamath, A., Jia, R., Liang, P., 2020. Selective question answering under domain shift. In: Proc.\ ACL. pp.\ 5684--5696. \url{https://doi.org/10.18653/v1/2020.acl-main.503}

\bibitem[Lample et~al.(2018)]{Lample2018}
Lample, G., Conneau, A., Denoyer, L., Ranzato, M., 2018. Unsupervised machine translation using monolingual corpora only. In: Proc.\ ICLR. arXiv:1711.00043. \url{https://doi.org/10.48550/arXiv.1711.00043}

\bibitem[Lewis et~al.(2020)]{Lewis2020}
Lewis, P., Perez, E., Piktus, A., Petroni, F., Karpukhin, V., Goyal, N., K\"uttler, H., Lewis, M., Yih, W., Rockt\"aschel, T., Riedel, S., Kiela, D., 2020. Retrieval-augmented generation for knowledge-intensive NLP tasks. In: Proc.\ NeurIPS. pp.\ 9459--9474. \url{https://doi.org/10.48550/arXiv.2005.11401}

\bibitem[Marcus and Maletic(2003)]{MarcusMaletic2003}
Marcus, A., Maletic, J.I., 2003. Recovering documentation-to-source-code traceability links using latent semantic indexing. In: Proc.\ 25th ICSE. pp.\ 349--357. \url{https://doi.org/10.1109/icse.2003.1201194}

\bibitem[Massey et~al.(2010)]{Massey2010}
Massey, A.K., Otto, P.N., Hayward, L.J., Ant\'on, A.I., 2010. Evaluating existing security and privacy requirements for legal compliance. Requir.\ Eng.\ 15~(1), 119--137. \url{https://doi.org/10.1007/s00766-009-0089-5}

\bibitem[McMillan et~al.(2009)]{McMillan2009}
McMillan, C., Poshyvanyk, D., Revelle, M., 2009. Combining textual and structural analysis of software artifacts for traceability link recovery. In: Proc.\ TEFSE Workshop. pp.\ 41--48. \url{https://doi.org/10.1109/tefse.2009.5069582}

\bibitem[NIST(2024)]{NISTIR8477}
National Institute of Standards and Technology, 2024. Mapping relationships between documentary standards, regulations, frameworks, and guidelines: developing cybersecurity and privacy concept mappings. NIST Internal Report (IR) 8477. \url{https://doi.org/10.6028/NIST.IR.8477}

\bibitem[Nogueira and Cho(2019)]{Nogueira2019}
Nogueira, R., Cho, K., 2019. Passage re-ranking with BERT. arXiv:1901.04085. \url{https://doi.org/10.48550/arXiv.1901.04085}

\bibitem[Oliveto et~al.(2010)]{Oliveto2010}
Oliveto, R., Gethers, M., Poshyvanyk, D., De~Lucia, A., 2010. On the equivalence of information retrieval methods for automated traceability link recovery. In: Proc.\ 18th IEEE ICPC. pp.\ 68--71. \url{https://doi.org/10.1109/icpc.2010.20}

\bibitem[Otto and Ant\'on(2007)]{OttoAnton2007}
Otto, P.N., Ant\'on, A.I., 2007. Addressing legal requirements in requirements engineering. In: Proc.\ 15th IEEE Int.\ Requirements Engineering Conf.\ (RE). pp.\ 5--14. \url{https://doi.org/10.1109/re.2007.65}

\bibitem[Recht et~al.(2019)]{Recht2019}
Recht, B., Roelofs, R., Schmidt, L., Shankar, V., 2019. Do ImageNet classifiers generalize to ImageNet? In: Proc.\ 36th Int.\ Conf.\ on Machine Learning (ICML). PMLR~97, pp.\ 5389--5400. \url{https://doi.org/10.48550/arXiv.1902.10811}

\bibitem[Reimers and Gurevych(2019)]{Reimers2019}
Reimers, N., Gurevych, I., 2019. Sentence-BERT: Sentence embeddings using Siamese BERT-networks. In: Proc.\ EMNLP. pp.\ 3982--3992. \url{https://doi.org/10.18653/v1/d19-1410}

\bibitem[Robertson and Zaragoza(2009)]{Robertson2009}
Robertson, S., Zaragoza, H., 2009. The probabilistic relevance framework: BM25 and beyond. Found.\ Trends Inf.\ Retr.\ 3~(4), 333--389. \url{https://doi.org/10.1561/1500000019}

\bibitem[Rocchio(1971)]{Rocchio1971}
Rocchio, J.J., 1971. Relevance feedback in information retrieval. In: Salton, G.\ (Ed.), The SMART Retrieval System: Experiments in Automatic Document Processing. Prentice Hall, pp.\ 313--323.

\bibitem[Schafer and Maxwell(2008)]{Maxwell2017}
Schafer, B., Maxwell, T., 2008. Concept and context in legal information retrieval. In: Proc.\ JURIX (Legal Knowledge and Information Systems). IOS Press, pp.\ 63--72. \url{https://doi.org/10.3233/978-1-58603-952-3-63}

\bibitem[Shao et~al.(2024)]{Shao2024}
Shao, Z., Gong, Y., Shen, Y., Huang, M., Duan, N., Chen, W., 2024. Assisting in writing Wikipedia-like articles from scratch with large language models. In: Proc.\ NAACL. pp.\ 6252--6278. \url{https://doi.org/10.18653/v1/2024.naacl-long.347}

\bibitem[Sim et~al.(2003)]{Sim2003}
Sim, S.E., Easterbrook, S., Holt, R.C., 2003. Using benchmarking to advance research: a challenge to software engineering. In: Proc.\ 25th Int.\ Conf.\ on Software Engineering (ICSE). pp.\ 74--83. \url{https://doi.org/10.1109/icse.2003.1201189}

\bibitem[Sleimi et~al.(2018)]{Sleimi2018}
Sleimi, A., Sannier, N., Sabetzadeh, M., Briand, L., Dann, J., 2018. Automated extraction of semantic legal metadata using natural language processing. In: Proc.\ 26th IEEE Int.\ Requirements Engineering Conf.\ (RE). pp.\ 124--135. \url{https://doi.org/10.1109/re.2018.00022}

\bibitem[Sun et~al.(2023)]{Sun2023}
Sun, W., Yan, L., Ma, X., Wang, S., Ren, P., Chen, Z., Yin, D., Ren, Z., 2023. Is ChatGPT good at search? Investigating large language models as re-ranking agents. In: Proc.\ EMNLP. pp.\ 14918--14937. \url{https://doi.org/10.18653/v1/2023.emnlp-main.923}

\bibitem[Thakur et~al.(2021)]{Thakur2021}
Thakur, N., Reimers, N., R\"uckl\'e, A., Srivastava, A., Gurevych, I., 2021. BEIR: A heterogeneous benchmark for zero-shot evaluation of information retrieval models. In: Proc.\ NeurIPS Datasets and Benchmarks Track. arXiv:2104.08663. \url{https://doi.org/10.48550/arXiv.2104.08663}

\bibitem[Voorhees(1994)]{Voorhees1994}
Voorhees, E.M., 1994. Query expansion using lexical-semantic relations. In: Proc.\ 17th ACM SIGIR. pp.\ 61--69. \url{https://doi.org/10.1007/978-1-4471-2099-5_7}

\bibitem[Voorhees(2000)]{Voorhees2000}
Voorhees, E.M., 2000. Variations in relevance judgments and the measurement of retrieval effectiveness. Inf.\ Process.\ Manage.\ 36~(5), 697--716. \url{https://doi.org/10.1016/s0306-4573(00)00010-8}

\bibitem[Yang et~al.(2019)]{Yang2019}
Yang, W., Lu, K., Yang, P., Lin, J., 2019. Critically examining the ``neural hype'': weak baselines and the additivity of effectiveness gains from neural ranking models. In: Proc.\ 42nd ACM SIGIR. pp.\ 1129--1132. \url{https://doi.org/10.1145/3331184.3331340}

\bibitem[Yang et~al.(2024)]{Yang2024}
Yang, J., Jimenez, C.E., Wettig, A., Liber, K., Narasimhan, K., Press, O., 2024. SWE-agent: Agent-computer interfaces enable automated software engineering. arXiv:2405.15793. \url{https://doi.org/10.52202/079017-1601}

\bibitem[Yao et~al.(2023)]{Yao2023}
Yao, S., Zhao, J., Yu, D., Du, N., Shafran, I., Narasimhan, K., Cao, Y., 2023. ReAct: Synergizing reasoning and acting in language models. In: Proc.\ ICLR. arXiv:2210.03629. \url{https://doi.org/10.48550/arXiv.2210.03629}

\bibitem[Zeni et~al.(2015)]{Zeni2015}
Zeni, N., Kiyavitskaya, N., Mich, L., Cordy, J.R., Mylopoulos, J., 2015. GaiusT: supporting the extraction of rights and obligations for regulatory compliance. Requir.\ Eng.\ 20~(1), 1--22. \url{https://doi.org/10.1007/s00766-013-0181-8}

\end{thebibliography}

\end{document}
